% homework/hw01/hw01.tex
% Year 7 - Homework Packet 1
% Covers Mathematics 1.1 (Adding & Subtracting Integers) and Science 1.1 (Cells).
%
% Styled after the Cambridge Lower Secondary Learner's Book, matching the lesson
% decks: teal for section headers, purple for activities, orange for anything the
% student must remember, blue for English help. The spine of both sections is
% ENGLISH -- every question that hides a language trap gets a word-help box, and
% the science answers are asked for in full sentences.
%
% Deliberate choices:
%   * No solid dark banners. Every header is a light tint with a coloured rule,
%     so the colour survives a cheap classroom printer without flooding it.
%   * No marks anywhere. This is practice, not an exam paper.
%   * Lato at 12pt with open leading -- this has to look like Year 7 homework,
%     not a university problem sheet.
%   * The source is pure ASCII: every typographic character goes through a
%     LaTeX command, so no editor or shell can mangle the encoding.
%
% Build:  pdflatex hw01.tex   (run twice)
% Teacher copy (adds the answer key):
%   pdflatex -jobname=hw01-teacher "\def\TEACHER{}% homework/hw01/hw01.tex
% Year 7 - Homework Packet 1
% Covers Mathematics 1.1 (Adding & Subtracting Integers) and Science 1.1 (Cells).
%
% Styled after the Cambridge Lower Secondary Learner's Book, matching the lesson
% decks: teal for section headers, purple for activities, orange for anything the
% student must remember, blue for English help. The spine of both sections is
% ENGLISH -- every question that hides a language trap gets a word-help box, and
% the science answers are asked for in full sentences.
%
% Deliberate choices:
%   * No solid dark banners. Every header is a light tint with a coloured rule,
%     so the colour survives a cheap classroom printer without flooding it.
%   * No marks anywhere. This is practice, not an exam paper.
%   * Lato at 12pt with open leading -- this has to look like Year 7 homework,
%     not a university problem sheet.
%   * The source is pure ASCII: every typographic character goes through a
%     LaTeX command, so no editor or shell can mangle the encoding.
%
% Build:  pdflatex hw01.tex   (run twice)
% Teacher copy (adds the answer key):
%   pdflatex -jobname=hw01-teacher "\def\TEACHER{}% homework/hw01/hw01.tex
% Year 7 - Homework Packet 1
% Covers Mathematics 1.1 (Adding & Subtracting Integers) and Science 1.1 (Cells).
%
% Styled after the Cambridge Lower Secondary Learner's Book, matching the lesson
% decks: teal for section headers, purple for activities, orange for anything the
% student must remember, blue for English help. The spine of both sections is
% ENGLISH -- every question that hides a language trap gets a word-help box, and
% the science answers are asked for in full sentences.
%
% Deliberate choices:
%   * No solid dark banners. Every header is a light tint with a coloured rule,
%     so the colour survives a cheap classroom printer without flooding it.
%   * No marks anywhere. This is practice, not an exam paper.
%   * Lato at 12pt with open leading -- this has to look like Year 7 homework,
%     not a university problem sheet.
%   * The source is pure ASCII: every typographic character goes through a
%     LaTeX command, so no editor or shell can mangle the encoding.
%
% Build:  pdflatex hw01.tex   (run twice)
% Teacher copy (adds the answer key):
%   pdflatex -jobname=hw01-teacher "\def\TEACHER{}% homework/hw01/hw01.tex
% Year 7 - Homework Packet 1
% Covers Mathematics 1.1 (Adding & Subtracting Integers) and Science 1.1 (Cells).
%
% Styled after the Cambridge Lower Secondary Learner's Book, matching the lesson
% decks: teal for section headers, purple for activities, orange for anything the
% student must remember, blue for English help. The spine of both sections is
% ENGLISH -- every question that hides a language trap gets a word-help box, and
% the science answers are asked for in full sentences.
%
% Deliberate choices:
%   * No solid dark banners. Every header is a light tint with a coloured rule,
%     so the colour survives a cheap classroom printer without flooding it.
%   * No marks anywhere. This is practice, not an exam paper.
%   * Lato at 12pt with open leading -- this has to look like Year 7 homework,
%     not a university problem sheet.
%   * The source is pure ASCII: every typographic character goes through a
%     LaTeX command, so no editor or shell can mangle the encoding.
%
% Build:  pdflatex hw01.tex   (run twice)
% Teacher copy (adds the answer key):
%   pdflatex -jobname=hw01-teacher "\def\TEACHER{}\input{hw01.tex}"

\documentclass[12pt,a4paper]{article}

% -- Answer key switch -------------------------------------------------------
\newif\ifkey
\ifdefined\TEACHER\keytrue\else
  \keyfalse        % \keyfalse = student copy   |   \keytrue = teacher copy
\fi

% -- House style -------------------------------------------------------------
% Shared by every packet. See homework/hw-style.tex.
\input{../hw-style.tex}

% -- Per-packet running head -------------------------------------------------
\fancyhead[L]{\small\color{black!55}Year 7\sep Homework 1}

\begin{document}

% ===========================================================================
% COVER BLOCK
% ===========================================================================
\thispagestyle{fancy}

\begin{tcolorbox}[enhanced, colback=cbpurple!7, colframe=cbpurple!7, arc=7pt,
                  borderline west={6pt}{0pt}{cbpurple},
                  left=18pt, right=14pt, top=14pt, bottom=14pt]
  {\color{cbpurple}\bfseries YEAR 7\sep UNIT 1\sep HOMEWORK 1}\\[6pt]
  {\color{cbpurple}\Huge\bfseries Integers and Cells}\\[10pt]
  {\color{black!70}Mathematics 1.1 --- Adding \& Subtracting Integers}\\
  {\color{black!70}Science 1.1 --- Cells and Microscopes}
\end{tcolorbox}

\vspace{10pt}
\noindent
\begin{tabularx}{\textwidth}{@{}X X@{}}
  \textbf{Name:} \blank[5.6cm] & \textbf{Class:} \blank[5.6cm] \\
  \textbf{Date given:} \blank[4.6cm] & \textbf{Date due:} \blank[5.0cm] \\
\end{tabularx}

\vspace{8pt}
\begin{activity}
  \begin{itemize}
    \item Write your answers \textbf{on this sheet}, in the spaces given. If you need
          more room, use the back of the page.
    \item In maths, \key{write the calculation before the answer}.
    \item In science, answer in \key{full English sentences}. ``Chloroplast'' is not a
          sentence; ``The green circles are chloroplasts'' is.
    \item If a word is new to you, look for the blue \eng{Word help} box on that page
          before you give up.
  \end{itemize}
\end{activity}

% ===========================================================================
% SECTION A -- MATHEMATICS
% ===========================================================================
\hwsection{Section A \textperiodcentered\ Mathematics 1.1 --- Adding and Subtracting Integers}

\begin{remember}
  An \key{integer} is a whole number that is positive, negative or zero --- never a
  fraction. On the number line, \key{negative} integers sit to the \key{left} of zero
  and \key{positive} integers to the \key{right}.

  \medskip
  \begin{tabularx}{\linewidth}{@{}lX@{}}
    Add a \key{positive}      & move \key{right} \\
    Add a \key{negative}      & move \key{left} \\
    Subtract a \key{positive} & move \key{left} \\
    Subtract a \key{negative} & move \key{right} --- because $a-(-b)=a+b$ \\
  \end{tabularx}

  \medskip
  Careful: only $-(-)$ becomes $+$. A single $+(-)$ still sends you left, so
  $5+(-3)=2$, \textbf{not} $8$.
\end{remember}

% -- A1 ---------------------------------------------------------------------
\qhead{A1}{Show it on the number line}

\noindent\textbf{(a)} Show $4 + (-3)$ on the number line below. Put a cross on the
number you \textbf{start} at, then draw an arrow for the move.

\vspace{12pt}
\noindent
\resizebox{0.92\textwidth}{!}{%
\begin{tikzpicture}[x=1cm,y=1cm]
  \draw[line width=1.1pt,cbteal,{Latex[length=3mm]}-{Latex[length=3mm]}] (-5.7,0) -- (8.7,0);
  \foreach \n in {-5,...,8}{
    \draw[line width=0.9pt,cbteal] (\n,-0.18) -- (\n,0.18);
    \node[below=4pt,font=\small] at (\n,-0.18) {$\n$};
  }
\end{tikzpicture}}

\vspace{10pt}
\noindent\hspace*{0.5em}$4 + (-3) = \blank[2.2cm]$

\vspace{12pt}
\noindent\textbf{(b)} Complete each sentence with \eng{higher than}, \eng{lower than},
\eng{warmer than} or \eng{colder than}.

\begin{enumerate}[label=(\roman*)]
  \item $-8$ is \blank[3.4cm] $-4$.
  \item A temperature of $3\dC$ is \blank[3.4cm] a temperature of $-9\dC$.
  \item Floor $-2$ of the car park is \blank[3.4cm] floor $1$.
\end{enumerate}

% -- A2 ---------------------------------------------------------------------
\qhead{A2}{Moving along the line}

\noindent Work these out. You may draw a small number line beside any of them.

\vspace{10pt}
\begin{multicols}{3}
\begin{enumerate}[label=(\alph*), itemsep=13pt]
  \item $-6+4 = \blank[1.6cm]$
  \item $5+(-9) = \blank[1.6cm]$
  \item $-3-5 = \blank[1.6cm]$
  \item $-2-(-7) = \blank[1.6cm]$
  \item $8+(-8) = \blank[1.6cm]$
  \item $-10-(-4) = \blank[1.6cm]$
  \item $0-(-6) = \blank[1.6cm]$
  \item $-7+15 = \blank[1.6cm]$
  \item $4-11 = \blank[1.6cm]$
  \item $-5-(-5) = \blank[1.6cm]$
  \item $6-8-(-5) = \blank[1.6cm]$
  \item $-9-6+(-1) = \blank[1.6cm]$
\end{enumerate}
\end{multicols}

% -- A3 ---------------------------------------------------------------------
\qhead{A3}{Say it, then write it}

\begin{wordhelp}[Word help --- these words decide the order of your calculation]
  \begin{tabularx}{\linewidth}{@{}lX@{}}
    \eng{subtract 5 from 8}   & you \textbf{start} at 8: $8-5$. The English gives the numbers in the \textbf{opposite order} to the calculation. \\
    \eng{take 6 away from 2}  & the same trap: $2-6$. \\
    \eng{7 less than 3}       & again the same trap: $3-7$. \\
    \eng{rises by / falls by} & the temperature goes up / the temperature goes down. \\
    \eng{deposit}             & to put money \textbf{into} an account $(+)$. \\
    \eng{withdraw}            & to take money \textbf{out of} an account $(-)$. \\
    \eng{owe} / \eng{a debt}  & money you must pay back. It makes your total \textbf{negative}. \\
    \eng{descend}             & to go \textbf{down}. To go up is to \eng{ascend}. \\
  \end{tabularx}
\end{wordhelp}

\noindent Turn each English sentence into a calculation, then work out the answer.
\key{Write the calculation first.}

\vspace{6pt}
\begin{enumerate}[label=\arabic*., itemsep=11pt]
  \item Subtract 9 from 4. \\[4pt]
        Calculation: \blank[6.2cm] \quad Answer: \blank[2.4cm]
  \item Take 6 away from $-2$. \\[4pt]
        Calculation: \blank[6.2cm] \quad Answer: \blank[2.4cm]
  \item What is 7 less than 3? \\[4pt]
        Calculation: \blank[6.2cm] \quad Answer: \blank[2.4cm]
  \item The temperature is $-8\dC$ and it rises by 5 degrees. \\[4pt]
        Calculation: \blank[6.2cm] \quad Answer: \blank[2.4cm]
  \item A submarine is at $-30$ m. It descends a further 12 m. \\[4pt]
        Calculation: \blank[6.2cm] \quad Answer: \blank[2.4cm]
  \item An account holds $-35$ dollars. The bank removes a debt of 20 dollars from it.
        \\[4pt]
        Calculation: \blank[6.2cm] \quad Answer: \blank[2.4cm]
\end{enumerate}

% -- A4 ---------------------------------------------------------------------
\qhead{A4}{Difference}

\begin{remember}[Remember: a difference is a gap]
  The \key{difference} between two numbers is how far apart they are on the line.
  Work it out with \key{bigger $-$ smaller}. A difference is \key{never negative}.
  \eng{How much warmer}, \eng{how much colder}, \eng{how much higher} and
  \eng{how many more} all ask for the same thing.
\end{remember}

\begin{enumerate}[label=\arabic*.]
  \item Find the difference between $-6$ and $2$. \hfill Answer: \blank[2.8cm]
  \item Find the difference between $-9$ and $-3$. \hfill Answer: \blank[2.8cm]
  \item How much warmer is $4\dC$ than $-13\dC$? Show your calculation.
        \wlines{1}
  \item These two questions look almost the same. They are not.
        \begin{enumerate}[label=(\alph*), itemsep=9pt]
          \item Find the difference between $-5$ and $6$. \hfill Answer: \blank[2.8cm]
          \item What is $-5$ minus $6$? \hfill Answer: \blank[2.8cm]
          \item In one full sentence, explain why the two answers are different.
                \wlines{2}
        \end{enumerate}
  \item One winter morning it is $18\dC$ in Ho Chi Minh City and $-24\dC$ in
        Ulaanbaatar. How much colder is Ulaanbaatar? Show your calculation.
        \wlines{1}
\end{enumerate}

% -- A5 ---------------------------------------------------------------------
\qhead{A5}{Word problems}

\noindent Read each one \key{all the way to the end} before you start writing.
Show the calculation, then the answer.

\begin{enumerate}[label=\textbf{\arabic*.}, itemsep=16pt]

  \item \textbf{The office lift.} Mr Bowen gets into the lift on floor 7. He goes
        \textbf{down} 11 floors, and then \textbf{up} 2 floors. Which floor is he on
        now?
        \begin{workspace}[2.8cm]\end{workspace}

  \item \textbf{A day in Sa Pa.} At 6 a.m. the temperature is $-3\dC$. By noon it has
        risen 11 degrees. By midnight it has fallen 14 degrees.
        \begin{enumerate}[label=(\alph*)]
          \item What is the temperature at midnight?
          \item What is the \textbf{difference} between the warmest and the coldest
                temperature of the day?
        \end{enumerate}
        \begin{workspace}[3.4cm]\end{workspace}

  \item \textbf{Six weeks of saving.} Mr Bowen's account is at $-40$ dollars. Every
        Monday he deposits 25 dollars. Every Friday he withdraws 25 dollars. How much is
        in the account after 6 weeks? Explain your answer in one full sentence.
        \begin{workspace}[3.1cm]\end{workspace}

\end{enumerate}

% -- A6 ---------------------------------------------------------------------
\qhead[cbgreen]{A6}{Now you write the question}

\begin{yourturn}[Your turn to be the teacher]
  So far the questions have given you the English, and you have written the maths.
  Now do it \textbf{backwards}. You are given the maths --- write the English.

  \medskip
  Your word problem must:
  \begin{itemize}
    \item be about something \textbf{real} --- money, temperature, floors in a
          building, or depth under the sea;
    \item use at least one word from the \eng{Word help} box on the earlier page,
          such as \eng{rises by}, \eng{falls by}, \eng{deposits}, \eng{withdraws},
          \eng{owes}, \eng{below} or \eng{descends};
    \item be written in \textbf{full English sentences}, and end with a question;
    \item have the answer that is already given to you.
  \end{itemize}
\end{yourturn}

\noindent\textbf{(a)} Write a word problem for \quad \expr{-5 + 8 = 3}
\wlines{4}

\vspace{14pt}
\noindent\textbf{(b)} Write a word problem for \quad \expr{12 - (-4) = 16}
\wlines{4}

% ===========================================================================
% SECTION B -- SCIENCE
% ===========================================================================
\hwsection{Section B \textperiodcentered\ Science 1.1 --- Cells and Microscopes}

\begin{remember}
  A \key{cell} is the smallest basic unit of all living organisms. An \key{organelle}
  is a tiny structure inside a cell that does one specific, important job to keep the
  cell alive. Every cell has a \key{cell membrane}, \key{cytoplasm}, a \key{nucleus}
  and \key{mitochondria}. A plant cell has all of those \key{plus} a \key{cell wall}
  made of \key{cellulose}, \key{chloroplasts} containing \key{chlorophyll}, and a
  \key{sap vacuole}.
\end{remember}

% -- B1 ---------------------------------------------------------------------
\qhead{B1}{Match the word to its job}

\noindent Write the correct letter in the box next to each word.

\vspace{8pt}
\noindent
\begin{tabularx}{\textwidth}{@{}p{4.2cm} @{\hspace{6pt}} c @{\hspace{12pt}} X@{}}
  \textbf{Word} & \textbf{Letter} & \textbf{Job} \\[2pt]
  1.\ Cell            & \tickbox & \textbf{A} \; The control centre of the cell --- the ``boss''. \\
  2.\ Organelle       & \tickbox & \textbf{B} \; A large space of sugars and water that keeps a plant cell firm. \\
  3.\ Nucleus         & \tickbox & \textbf{C} \; Where energy is released from food. \\
  4.\ Mitochondria    & \tickbox & \textbf{D} \; The smallest basic unit of all living organisms. \\
  5.\ Cell wall       & \tickbox & \textbf{E} \; A tool that uses lenses to bend light and magnify an image. \\
  6.\ Chloroplast     & \tickbox & \textbf{F} \; A tiny structure inside a cell that does one specific job. \\
  7.\ Sap vacuole     & \tickbox & \textbf{G} \; A strong, stiff outer layer that holds a plant cell in shape. \\
  8.\ Microscope      & \tickbox & \textbf{H} \; Where a plant makes its food, using sunlight. \\
\end{tabularx}

% -- B2 ---------------------------------------------------------------------
\qhead{B2}{Animal, plant, or both?}

\noindent Put a tick ($\checkmark$) in \textbf{one} column for each organelle.

\vspace{8pt}
\noindent
\begin{tabularx}{\textwidth}{|X|>{\centering\arraybackslash}p{3.1cm}|>{\centering\arraybackslash}p{2.9cm}|>{\centering\arraybackslash}p{1.9cm}|}
  \hline
  \textbf{Organelle} & \textbf{Animal cells only} & \textbf{Plant cells only} & \textbf{Both} \\
  \hline
  Cell membrane & & & \\ \hline
  Cytoplasm     & & & \\ \hline
  Nucleus       & & & \\ \hline
  Mitochondria  & & & \\ \hline
  Cell wall     & & & \\ \hline
  Chloroplast   & & & \\ \hline
  Sap vacuole   & & & \\ \hline
\end{tabularx}

% -- B4 ---------------------------------------------------------------------
\qhead{B3}{Answer in full sentences}

\begin{wordhelp}[Sentence starters --- you may copy these]
  ``A plant cell needs a cell wall because\ldots''\\
  ``One difference is that\ldots\ Another difference is that\ldots''\\
  ``A microscope is a tool that\ldots''
\end{wordhelp}

\begin{enumerate}[label=\arabic*., itemsep=14pt]

  \item An animal has a skeleton to hold it up. A plant does not. Explain why every
        plant cell needs a cell wall.
        \wlines{2}

  \item Give \textbf{two} ways a cell wall is different from a cell membrane.
        Write each one as a full sentence.
        \wlines{4}

  \item Under the microscope, a leaf cell is full of small green circles. Name them,
        and name the green substance inside them.
        \wlines{2}

  \item What does a microscope do? Use the words \key{lens} and \key{magnify} in your
        answer.
        \wlines{3}

\end{enumerate}

% -- B5 ---------------------------------------------------------------------
\qhead{B4}{A model of a cell}

\noindent Mr Bowen builds a model of an animal cell in his kitchen. He uses a clear
plastic bag, filled with jelly, with a plum and four red beans pushed into the middle
of it.

\vspace{8pt}
\noindent\textbf{(a)} Which organelle does each item represent?

\vspace{6pt}
\noindent
\begin{tabularx}{\textwidth}{|p{5cm}|X|}
  \hline
  \textbf{Item in the model} & \textbf{Organelle it represents} \\
  \hline
  The plastic bag \rule{0pt}{1.9em} & \\ \hline
  The jelly       \rule{0pt}{1.9em} & \\ \hline
  The plum        \rule{0pt}{1.9em} & \\ \hline
  The four beans  \rule{0pt}{1.9em} & \\ \hline
\end{tabularx}

\vspace{14pt}
\noindent\textbf{(b)} Name \textbf{one} thing he could add to turn it into a model of a
\textbf{plant} cell, and say which organelle it would represent.
\wlines{2}

% -- B3 ---------------------------------------------------------------------
\qhead{B5}{Label the plant cell}

\noindent Write the name of each numbered part on the lines below the diagram.

\vspace{12pt}
\noindent
\resizebox{\textwidth}{!}{%
\begin{tikzpicture}[x=1cm,y=1cm,
    lbl/.style={circle, draw=cbteal, fill=white, line width=1.1pt,
                inner sep=1.8pt, minimum size=8mm, font=\bfseries}]

  % -- the cell itself. The wall band is deliberately 0.7 wide so that the
  %    wall (1) and the membrane (3) are visibly two different things.
  \fill[cbgreen!35, rounded corners=12pt] (0,0) rectangle (11,6.4);        % cell wall
  \fill[cbteal!10, rounded corners=8pt]  (0.7,0.7) rectangle (10.3,5.7);   % cytoplasm
  \draw[line width=1.2pt, black!60, rounded corners=8pt] (0.7,0.7) rectangle (10.3,5.7);
  \draw[line width=1.6pt, cbgreen!70!black, rounded corners=12pt] (0,0) rectangle (11,6.4);

  % -- sap vacuole --
  \fill[cbblue!18] (5.5,2.95) ellipse (3.2 and 1.65);
  \draw[line width=0.9pt, cbblue!60] (5.5,2.95) ellipse (3.2 and 1.65);

  % -- nucleus --
  \fill[cbpurple!30] (2.0,4.35) circle (0.85);
  \draw[line width=0.9pt, cbpurple!80] (2.0,4.35) circle (0.85);
  \fill[cbpurple!60] (2.0,4.35) circle (0.28);

  % -- chloroplasts --
  \foreach \x/\y/\r in {1.7/1.5/20, 3.2/1.15/-15, 5.0/1.05/10,
                        6.8/1.15/-20, 8.5/1.5/15, 9.6/2.6/-10, 9.6/4.2/25}{
    \begin{scope}[shift={(\x,\y)}, rotate=\r]
      \fill[cbgreen!75] (0,0) ellipse (0.46 and 0.24);
      \draw[line width=0.7pt, cbgreen!45!black] (0,0) ellipse (0.46 and 0.24);
    \end{scope}}

  % -- mitochondria. Kept away from x = 5.5 so the cytoplasm leader (2) has a
  %    clear, empty patch of cytoplasm to land on.
  \foreach \x/\y/\r in {3.3/5.1/-12, 8.4/5.1/16, 1.35/2.7/70}{
    \begin{scope}[shift={(\x,\y)}, rotate=\r]
      \fill[cborange!65] (0,0) ellipse (0.42 and 0.2);
      \draw[line width=0.7pt, cborange!80!black] (0,0) ellipse (0.42 and 0.2);
      \draw[line width=0.5pt, cborange!80!black] (-0.25,0.05) -- (-0.05,-0.06) -- (0.15,0.06);
    \end{scope}}

  % -- leader lines and numbers. Each one ends on a dot so there is no doubt
  %    about which structure it is pointing at.
  \foreach \sx/\sy/\ex/\ey/\lx/\ly/\n in {%
      2.2/6.4/2.2/7.55/2.2/7.95/1,          % cell wall  (the green band)
      5.5/5.35/7.9/7.55/8.0/7.95/2,         % cytoplasm  (empty patch)
      0.7/4.9/-1.5/6.25/-1.9/6.4/3,         % cell membrane (the black line)
      2.0/4.35/-1.5/4.05/-1.9/4.0/4,        % nucleus
      5.5/2.95/12.6/3.3/13.0/3.35/5,        % sap vacuole
      1.7/1.5/-1.5/1.1/-1.9/1.05/6,         % chloroplast
      8.4/5.1/12.6/5.5/13.0/5.55/7}{        % mitochondrion
    \draw[line width=0.7pt, black!65] (\sx,\sy) -- (\ex,\ey);
    \fill[black!65] (\sx,\sy) circle (0.075);
    \node[lbl] at (\lx,\ly) {\n};}

  % keep the bounding box sensible
  \path (-2.6,-0.5) rectangle (13.9,8.6);
\end{tikzpicture}}

\vspace{16pt}

% A tabular rather than multicols, so the seven answer lines can never be split
% across a page break away from their diagram.
\noindent
\begin{tabular}{@{}p{0.48\textwidth} p{0.48\textwidth}@{}}
  1.\ \blank[5.4cm] & 5.\ \blank[5.4cm] \\[30pt]
  2.\ \blank[5.4cm] & 6.\ \blank[5.4cm] \\[30pt]
  3.\ \blank[5.4cm] & 7.\ \blank[5.4cm] \\[30pt]
  4.\ \blank[5.4cm] &                   \\
\end{tabular}

\vspace{10pt}
\noindent{\small\color{black!60}Hint: number 3 is the very thin layer just inside the
outer green edge. Numbers 1 and 3 are easy to confuse --- one is thick and stiff, one is
thin and flexible.}

% ===========================================================================
% ANSWER KEY (teacher copy only)
% ===========================================================================
\ifkey
\newpage
\hwsection[cbred]{Answer Key --- Teacher Copy}

% The key is for an adult reading it at a desk, not a 12-year-old writing on it,
% so tighten everything back up.
\linespread{1.0}\selectfont
\setlist[enumerate]{itemsep=2pt, topsep=3pt, leftmargin=*}
\setlist[itemize]{itemsep=2pt, topsep=3pt, leftmargin=*}

\qhead[cbred]{A1}{Show it on the number line}
\noindent (a) Cross on $4$, arrow three steps \emph{left}, landing on $1$. So
$4+(-3)=1$. The cross must be on 4, not on 3 --- adding a negative means starting at
the first number and moving left.\\
(b) (i) lower than \quad (ii) warmer than \quad (iii) lower than.
{\small\ (``higher than'' is also acceptable in (iii) if the sentence is reversed ---
mark the reasoning, not the word.)}

\qhead[cbred]{A2}{Moving along the line}
\noindent
(a) $-2$ \quad (b) $-4$ \quad (c) $-8$ \quad (d) $5$ \quad (e) $0$ \quad (f) $-6$ \\
(g) $6$ \quad (h) $8$ \quad (i) $-7$ \quad (j) $0$ \quad (k) $3$ \quad (l) $-16$

\smallskip
\noindent{\small Common errors: (d) and (f) --- students flip \emph{any} pair of signs
to $+$. Only $-(-)$ becomes $+$. (b) $5+(-9)$ is a single $+(-)$, so it still goes
left.}

\qhead[cbred]{A3}{Say it, then write it}
\begin{enumerate}[label=\arabic*., itemsep=3pt]
  \item $4-9=-5$
  \item $-2-6=-8$
  \item $3-7=-4$
  \item $-8+5=-3$
  \item $-30-12=-42$
  \item $-35-(-20)=-15$
\end{enumerate}
\noindent{\small Questions 1--3 are the ``reversed order'' trap; 6 is the ``removing a
debt adds'' rule --- the only place it is tested now, so it is worth checking. If a student writes $9-4$, $6-(-2)$ or $7-3$, that is the English
error, not an arithmetic one --- say so, and make them re-read the sentence.}

\qhead[cbred]{A4}{Difference}
\begin{enumerate}[label=\arabic*., itemsep=3pt]
  \item $8$
  \item $6$
  \item $4-(-13)=17$ degrees
  \item (a) $11$ \; (b) $-11$ \; (c) A difference is the gap between the numbers so it
        is never negative, but $-5-6$ is an instruction to move 6 further left.
  \item $18-(-24)=42$ degrees
\end{enumerate}

\qhead[cbred]{A5}{Word problems}
\begin{enumerate}[label=\arabic*., itemsep=3pt]
  \item $7-11+2=-2$. Floor $-2$ (B2).
  \item (a) $-3+11-14=-6\dC$. \; (b) Warmest $8\dC$ at noon, coldest $-6\dC$:
        $8-(-6)=14$ degrees.
  \item Still $-40$ dollars. Each week $+25-25=0$, so six weeks change nothing. This one
        rewards reading to the end before calculating.
\end{enumerate}

\qhead[cbred]{A6}{Now you write the question}
\noindent There is no single right answer --- check the four conditions in the green
box, in this order:
\begin{enumerate}[label=\arabic*., itemsep=3pt]
  \item \textbf{Does it start in the right place?} For (a) the story must \emph{begin}
        at $-5$, not arrive there. This is the same ``which number comes first''
        problem as A3, and it is the one most of them will get wrong.
  \item \textbf{Does the signal word match the sign?} (a) needs an \emph{up} word
        (rises, deposits, climbs) for the $+8$. (b) needs a \emph{removing a negative}
        --- a debt cancelled, a weight taken off --- not simply ``adds 4''.
  \item \textbf{Full sentences, ending in a question?}
  \item \textbf{Does their own answer come out as 3, and as 16?}
\end{enumerate}

\smallskip
\noindent{\small Two that work: (a) ``The temperature in Sa Pa was $-5\dC$ at
midnight. By noon it had risen by 8 degrees. What was the temperature at noon?''
(b) ``Mr Bowen had 12 dollars, and he also owed his brother 4 dollars. His brother told
him to forget the debt. How much is Mr Bowen worth now?'' \\
Part (b) is the harder one by a distance --- expect most of the class to write a plain
addition. That is worth a whole-class discussion rather than a red pen: ask them what
in real life makes you \emph{richer} when it is taken away.}

\qhead[cbred]{B1}{Match the word to its job}
\noindent 1--D \quad 2--F \quad 3--A \quad 4--C \quad 5--G \quad 6--H \quad 7--B \quad 8--E

\qhead[cbred]{B2}{Animal, plant, or both?}
\noindent Both: cell membrane, cytoplasm, nucleus, mitochondria. \\
Plant only: cell wall, chloroplast, sap vacuole. \\
{\small Nothing is animal-only. Expect some students to tick that column anyway ---
it is worth saying out loud that a plant cell has everything an animal cell has,
\emph{plus} extras.}

\qhead[cbred]{B3}{Answer in full sentences}
\begin{enumerate}[label=\arabic*., itemsep=3pt]
  \item A plant has no skeleton, so each cell needs a stiff wall around it to hold its
        shape and hold the plant up.
  \item Any two: the wall is thick and stiff while the membrane is thin and flexible;
        the wall is made of cellulose; the wall is only in plant cells while every cell
        has a membrane; the membrane controls what enters and leaves, the wall does not.
  \item They are chloroplasts, and the green substance inside them is chlorophyll.
  \item A microscope is a tool that uses curved pieces of glass called lenses to bend
        light and magnify an image, making it look much bigger than it really is.
\end{enumerate}

\qhead[cbred]{B4}{A model of a cell}
\noindent (a) bag --- cell membrane; jelly --- cytoplasm; plum --- nucleus;
beans --- mitochondria.

\smallskip
\noindent (b) Any one, with its organelle named: a shoebox or rigid container around the
bag --- cell wall; a water balloon in the middle --- sap vacuole; green sweets ---
chloroplasts.

\qhead[cbred]{B5}{Label the plant cell}
\noindent 1 --- cell wall \quad 2 --- cytoplasm \quad 3 --- cell membrane \quad
4 --- nucleus \quad 5 --- sap vacuole \quad 6 --- chloroplast \quad
7 --- mitochondrion

\fi

\end{document}
"

\documentclass[12pt,a4paper]{article}

% -- Answer key switch -------------------------------------------------------
\newif\ifkey
\ifdefined\TEACHER\keytrue\else
  \keyfalse        % \keyfalse = student copy   |   \keytrue = teacher copy
\fi

% -- House style -------------------------------------------------------------
% Shared by every packet. See homework/hw-style.tex.
% homework/hw-style.tex
% Shared house style for every Year 7 homework packet. \input this from the
% packet file, right after \documentclass and the answer-key switch.
%
% Do not fork this file per packet. If a packet needs a different look, the
% question is whether every packet should have it.
%
% The choices here are deliberate and were arrived at by printing and looking:
%   * No solid dark banners. Every header is a light tint with a coloured spine
%     on the left, so colour reads clearly without flooding a school printer.
%   * No marks and no mark totals. These are practice, not exam papers.
%   * Lato at 12pt with open leading, plus mathastext so inline maths is Lato
%     too. Without mathastext every "-6 + 4" is Computer Modern serif and the
%     sheet looks half-typeset.
%   * Writing lines are spaced for Year 7 handwriting, not adult margin notes.
%
% Provides: \wlines \blank \tickbox \sep \qmk-free marks (none) \hwsection
% \qhead \expr, the workspace box, and the remember / wordhelp / activity /
% yourturn coloured boxes. Palette matches the lesson decks exactly.

% -- Packages ----------------------------------------------------------------
\usepackage[T1]{fontenc}
\usepackage[utf8]{inputenc}
\usepackage[default]{lato}
\usepackage[a4paper,top=1.7cm,bottom=1.7cm,left=2.0cm,right=2.0cm,headheight=15pt]{geometry}
\usepackage{amsmath,amssymb}
\usepackage{xcolor}
\usepackage[most]{tcolorbox}
\usepackage{enumitem}
\usepackage{tikz}
\usetikzlibrary{arrows.meta}
\usepackage{array}
\usepackage{tabularx}
\usepackage{fancyhdr}
\usepackage{multicol}
\usepackage{needspace}
\usepackage{graphicx}

% Make maths use Lato too. Without this, every "-6 + 4" on the page is set in
% Computer Modern serif and the sheet looks half-typeset. Must come last.
\usepackage[italic]{mathastext}

\linespread{1.06}\selectfont

% -- The Learner's Book palette (same hex values as the lesson decks) ---------
\definecolor{cbteal}{HTML}{0087A8}
\definecolor{cbpurple}{HTML}{5C2483}
\definecolor{cborange}{HTML}{C25E12}
\definecolor{cbgreen}{HTML}{4A8B23}
\definecolor{cbred}{HTML}{C8102E}
\definecolor{cbblue}{HTML}{1A5FA8}
\definecolor{rulegray}{HTML}{9AA5AD}

% -- Answer lines and blanks -------------------------------------------------
% \wlines{n} draws n ruled writing lines, spaced wide enough for Year 7
% handwriting rather than for an adult with a fine pen.
\makeatletter
\newcommand{\wlines}[1]{%
  \par\vspace{0.75em}%
  \@tempcnta=0\relax
  \@whilenum\@tempcnta<#1\do{%
    \hbox to \linewidth{\color{rulegray}\leaders\hrule height 0.5pt\hfill}%
    \vspace{1.5em}%
    \advance\@tempcnta by 1\relax}%
  \par\vspace{-0.8em}%
}
\makeatother

% \blank[width] is an inline gap to write a short answer in.
\newcommand{\blank}[1][2.6cm]{\underline{\hspace{#1}}}

% A boxed space for showing working.
\newtcolorbox{workspace}[1][3.4cm]{%
  colback=white, colframe=rulegray, boxrule=0.5pt, arc=5pt,
  height=#1, left=6pt, right=6pt, top=4pt, bottom=4pt,
  before skip=9pt, after skip=11pt}

% An empty tick box.
\newcommand{\tickbox}{\fbox{\rule{0pt}{1.15em}\rule{1.15em}{0pt}}}

% A middot separator.
\newcommand{\sep}{\,\textperiodcentered\,}

% -- Section and question headings -------------------------------------------
% Light tint plus a fat coloured spine on the left. No solid dark banner: the
% colour reads clearly but barely costs the printer anything.
% needspace stops a heading stranding alone at the foot of a page. Kept modest:
% too large and it pushes headings over, leaving a worse hole than it prevents.
\newcommand{\hwsection}[2][cbteal]{%
  \needspace{8\baselineskip}%
  \par\vspace{13pt}%
  \noindent
  \begin{tcolorbox}[enhanced, nobeforeafter, width=\textwidth,
      colback=#1!7, colframe=#1!7, arc=5pt,
      borderline west={5pt}{0pt}{#1},
      left=14pt, right=10pt, top=7pt, bottom=7pt]
    {\large\bfseries\color{#1}#2}
  \end{tcolorbox}%
  \par\vspace{9pt}}

% \qhead[colour]{A2}{Moving along the line} -- a question heading.
\newcommand{\qhead}[3][cbteal]{%
  \needspace{6\baselineskip}%
  \par\vspace{11pt}%
  \noindent{\large\bfseries\color{#1}#2}\quad{\large\bfseries #3}\par\vspace{4pt}}

% -- Coloured boxes ----------------------------------------------------------
% Shared look: rounded, light fill, coloured rule, coloured title text on a
% pale title strip rather than a dark one.
\tcbset{hwbox/.style={
  enhanced, breakable, arc=6pt, boxrule=1pt,
  left=11pt, right=11pt, top=8pt, bottom=8pt,
  fonttitle=\bfseries, attach boxed title to top left={xshift=12pt, yshift=-8pt},
  boxed title style={arc=4pt, boxrule=1pt, left=7pt, right=7pt, top=2pt, bottom=2pt},
  top=13pt}}

% Orange: things to remember. Matches the deck's "Write This Down".
\newtcolorbox{remember}[1][Remember from the lesson]{hwbox,
  colback=cborange!6, colframe=cborange, coltitle=cborange,
  colbacktitle=white, title={#1}}

% Blue: English help. The barrier in this class is the wording, not the maths.
\newtcolorbox{wordhelp}[1][Word help]{hwbox,
  colback=cbblue!6, colframe=cbblue, coltitle=cbblue,
  colbacktitle=white, title={#1}}

% Purple: an activity or instruction.
\newtcolorbox{activity}[1][How to do this homework]{hwbox,
  colback=cbpurple!6, colframe=cbpurple, coltitle=cbpurple,
  colbacktitle=white, title={#1}}

% Green: your turn / a friendly nudge.
\newtcolorbox{yourturn}[1][Your turn]{hwbox,
  colback=cbgreen!7, colframe=cbgreen, coltitle=cbgreen!75!black,
  colbacktitle=white, title={#1}}

% -- Shorthands --------------------------------------------------------------
\newcommand{\dC}{\,^{\circ}\mathrm{C}}          % use inside maths: $-8\dC$
\newcommand{\key}[1]{{\color{cborange}\textbf{#1}}}
\newcommand{\eng}[1]{{\color{cbblue}\textbf{#1}}}

\setlist[enumerate]{itemsep=8pt, topsep=5pt, leftmargin=*}
\setlist[itemize]{itemsep=4pt, topsep=4pt, leftmargin=*}
\renewcommand{\arraystretch}{1.45}

% -- An expression the student is asked to work from -------------------------
\newcommand{\expr}[1]{%
  \tcbox[on line, colback=cbgreen!10, colframe=cbgreen!55, boxrule=1pt, arc=4pt,
         left=9pt, right=9pt, top=4pt, bottom=4pt]{\large\bfseries $#1$}}

% -- Running head and foot ---------------------------------------------------
% The packet overrides \fancyhead[L] with its own title.
\pagestyle{fancy}
\fancyhf{}
\fancyhead[L]{\small\color{black!55}Year 7\sep Homework}
\fancyhead[R]{\small\color{black!55}Name: \blank[3.4cm]}
\fancyfoot[C]{\small\color{black!55}Page \thepage}
\renewcommand{\headrulewidth}{0pt}
\renewcommand{\headrule}{\vspace{2pt}\hbox to\headwidth{\color{cbteal!45}\leaders\hrule height 1.2pt\hfill}}



% -- Per-packet running head -------------------------------------------------
\fancyhead[L]{\small\color{black!55}Year 7\sep Homework 1}

\begin{document}

% ===========================================================================
% COVER BLOCK
% ===========================================================================
\thispagestyle{fancy}

\begin{tcolorbox}[enhanced, colback=cbpurple!7, colframe=cbpurple!7, arc=7pt,
                  borderline west={6pt}{0pt}{cbpurple},
                  left=18pt, right=14pt, top=14pt, bottom=14pt]
  {\color{cbpurple}\bfseries YEAR 7\sep UNIT 1\sep HOMEWORK 1}\\[6pt]
  {\color{cbpurple}\Huge\bfseries Integers and Cells}\\[10pt]
  {\color{black!70}Mathematics 1.1 --- Adding \& Subtracting Integers}\\
  {\color{black!70}Science 1.1 --- Cells and Microscopes}
\end{tcolorbox}

\vspace{10pt}
\noindent
\begin{tabularx}{\textwidth}{@{}X X@{}}
  \textbf{Name:} \blank[5.6cm] & \textbf{Class:} \blank[5.6cm] \\
  \textbf{Date given:} \blank[4.6cm] & \textbf{Date due:} \blank[5.0cm] \\
\end{tabularx}

\vspace{8pt}
\begin{activity}
  \begin{itemize}
    \item Write your answers \textbf{on this sheet}, in the spaces given. If you need
          more room, use the back of the page.
    \item In maths, \key{write the calculation before the answer}.
    \item In science, answer in \key{full English sentences}. ``Chloroplast'' is not a
          sentence; ``The green circles are chloroplasts'' is.
    \item If a word is new to you, look for the blue \eng{Word help} box on that page
          before you give up.
  \end{itemize}
\end{activity}

% ===========================================================================
% SECTION A -- MATHEMATICS
% ===========================================================================
\hwsection{Section A \textperiodcentered\ Mathematics 1.1 --- Adding and Subtracting Integers}

\begin{remember}
  An \key{integer} is a whole number that is positive, negative or zero --- never a
  fraction. On the number line, \key{negative} integers sit to the \key{left} of zero
  and \key{positive} integers to the \key{right}.

  \medskip
  \begin{tabularx}{\linewidth}{@{}lX@{}}
    Add a \key{positive}      & move \key{right} \\
    Add a \key{negative}      & move \key{left} \\
    Subtract a \key{positive} & move \key{left} \\
    Subtract a \key{negative} & move \key{right} --- because $a-(-b)=a+b$ \\
  \end{tabularx}

  \medskip
  Careful: only $-(-)$ becomes $+$. A single $+(-)$ still sends you left, so
  $5+(-3)=2$, \textbf{not} $8$.
\end{remember}

% -- A1 ---------------------------------------------------------------------
\qhead{A1}{Show it on the number line}

\noindent\textbf{(a)} Show $4 + (-3)$ on the number line below. Put a cross on the
number you \textbf{start} at, then draw an arrow for the move.

\vspace{12pt}
\noindent
\resizebox{0.92\textwidth}{!}{%
\begin{tikzpicture}[x=1cm,y=1cm]
  \draw[line width=1.1pt,cbteal,{Latex[length=3mm]}-{Latex[length=3mm]}] (-5.7,0) -- (8.7,0);
  \foreach \n in {-5,...,8}{
    \draw[line width=0.9pt,cbteal] (\n,-0.18) -- (\n,0.18);
    \node[below=4pt,font=\small] at (\n,-0.18) {$\n$};
  }
\end{tikzpicture}}

\vspace{10pt}
\noindent\hspace*{0.5em}$4 + (-3) = \blank[2.2cm]$

\vspace{12pt}
\noindent\textbf{(b)} Complete each sentence with \eng{higher than}, \eng{lower than},
\eng{warmer than} or \eng{colder than}.

\begin{enumerate}[label=(\roman*)]
  \item $-8$ is \blank[3.4cm] $-4$.
  \item A temperature of $3\dC$ is \blank[3.4cm] a temperature of $-9\dC$.
  \item Floor $-2$ of the car park is \blank[3.4cm] floor $1$.
\end{enumerate}

% -- A2 ---------------------------------------------------------------------
\qhead{A2}{Moving along the line}

\noindent Work these out. You may draw a small number line beside any of them.

\vspace{10pt}
\begin{multicols}{3}
\begin{enumerate}[label=(\alph*), itemsep=13pt]
  \item $-6+4 = \blank[1.6cm]$
  \item $5+(-9) = \blank[1.6cm]$
  \item $-3-5 = \blank[1.6cm]$
  \item $-2-(-7) = \blank[1.6cm]$
  \item $8+(-8) = \blank[1.6cm]$
  \item $-10-(-4) = \blank[1.6cm]$
  \item $0-(-6) = \blank[1.6cm]$
  \item $-7+15 = \blank[1.6cm]$
  \item $4-11 = \blank[1.6cm]$
  \item $-5-(-5) = \blank[1.6cm]$
  \item $6-8-(-5) = \blank[1.6cm]$
  \item $-9-6+(-1) = \blank[1.6cm]$
\end{enumerate}
\end{multicols}

% -- A3 ---------------------------------------------------------------------
\qhead{A3}{Say it, then write it}

\begin{wordhelp}[Word help --- these words decide the order of your calculation]
  \begin{tabularx}{\linewidth}{@{}lX@{}}
    \eng{subtract 5 from 8}   & you \textbf{start} at 8: $8-5$. The English gives the numbers in the \textbf{opposite order} to the calculation. \\
    \eng{take 6 away from 2}  & the same trap: $2-6$. \\
    \eng{7 less than 3}       & again the same trap: $3-7$. \\
    \eng{rises by / falls by} & the temperature goes up / the temperature goes down. \\
    \eng{deposit}             & to put money \textbf{into} an account $(+)$. \\
    \eng{withdraw}            & to take money \textbf{out of} an account $(-)$. \\
    \eng{owe} / \eng{a debt}  & money you must pay back. It makes your total \textbf{negative}. \\
    \eng{descend}             & to go \textbf{down}. To go up is to \eng{ascend}. \\
  \end{tabularx}
\end{wordhelp}

\noindent Turn each English sentence into a calculation, then work out the answer.
\key{Write the calculation first.}

\vspace{6pt}
\begin{enumerate}[label=\arabic*., itemsep=11pt]
  \item Subtract 9 from 4. \\[4pt]
        Calculation: \blank[6.2cm] \quad Answer: \blank[2.4cm]
  \item Take 6 away from $-2$. \\[4pt]
        Calculation: \blank[6.2cm] \quad Answer: \blank[2.4cm]
  \item What is 7 less than 3? \\[4pt]
        Calculation: \blank[6.2cm] \quad Answer: \blank[2.4cm]
  \item The temperature is $-8\dC$ and it rises by 5 degrees. \\[4pt]
        Calculation: \blank[6.2cm] \quad Answer: \blank[2.4cm]
  \item A submarine is at $-30$ m. It descends a further 12 m. \\[4pt]
        Calculation: \blank[6.2cm] \quad Answer: \blank[2.4cm]
  \item An account holds $-35$ dollars. The bank removes a debt of 20 dollars from it.
        \\[4pt]
        Calculation: \blank[6.2cm] \quad Answer: \blank[2.4cm]
\end{enumerate}

% -- A4 ---------------------------------------------------------------------
\qhead{A4}{Difference}

\begin{remember}[Remember: a difference is a gap]
  The \key{difference} between two numbers is how far apart they are on the line.
  Work it out with \key{bigger $-$ smaller}. A difference is \key{never negative}.
  \eng{How much warmer}, \eng{how much colder}, \eng{how much higher} and
  \eng{how many more} all ask for the same thing.
\end{remember}

\begin{enumerate}[label=\arabic*.]
  \item Find the difference between $-6$ and $2$. \hfill Answer: \blank[2.8cm]
  \item Find the difference between $-9$ and $-3$. \hfill Answer: \blank[2.8cm]
  \item How much warmer is $4\dC$ than $-13\dC$? Show your calculation.
        \wlines{1}
  \item These two questions look almost the same. They are not.
        \begin{enumerate}[label=(\alph*), itemsep=9pt]
          \item Find the difference between $-5$ and $6$. \hfill Answer: \blank[2.8cm]
          \item What is $-5$ minus $6$? \hfill Answer: \blank[2.8cm]
          \item In one full sentence, explain why the two answers are different.
                \wlines{2}
        \end{enumerate}
  \item One winter morning it is $18\dC$ in Ho Chi Minh City and $-24\dC$ in
        Ulaanbaatar. How much colder is Ulaanbaatar? Show your calculation.
        \wlines{1}
\end{enumerate}

% -- A5 ---------------------------------------------------------------------
\qhead{A5}{Word problems}

\noindent Read each one \key{all the way to the end} before you start writing.
Show the calculation, then the answer.

\begin{enumerate}[label=\textbf{\arabic*.}, itemsep=16pt]

  \item \textbf{The office lift.} Mr Bowen gets into the lift on floor 7. He goes
        \textbf{down} 11 floors, and then \textbf{up} 2 floors. Which floor is he on
        now?
        \begin{workspace}[2.8cm]\end{workspace}

  \item \textbf{A day in Sa Pa.} At 6 a.m. the temperature is $-3\dC$. By noon it has
        risen 11 degrees. By midnight it has fallen 14 degrees.
        \begin{enumerate}[label=(\alph*)]
          \item What is the temperature at midnight?
          \item What is the \textbf{difference} between the warmest and the coldest
                temperature of the day?
        \end{enumerate}
        \begin{workspace}[3.4cm]\end{workspace}

  \item \textbf{Six weeks of saving.} Mr Bowen's account is at $-40$ dollars. Every
        Monday he deposits 25 dollars. Every Friday he withdraws 25 dollars. How much is
        in the account after 6 weeks? Explain your answer in one full sentence.
        \begin{workspace}[3.1cm]\end{workspace}

\end{enumerate}

% -- A6 ---------------------------------------------------------------------
\qhead[cbgreen]{A6}{Now you write the question}

\begin{yourturn}[Your turn to be the teacher]
  So far the questions have given you the English, and you have written the maths.
  Now do it \textbf{backwards}. You are given the maths --- write the English.

  \medskip
  Your word problem must:
  \begin{itemize}
    \item be about something \textbf{real} --- money, temperature, floors in a
          building, or depth under the sea;
    \item use at least one word from the \eng{Word help} box on the earlier page,
          such as \eng{rises by}, \eng{falls by}, \eng{deposits}, \eng{withdraws},
          \eng{owes}, \eng{below} or \eng{descends};
    \item be written in \textbf{full English sentences}, and end with a question;
    \item have the answer that is already given to you.
  \end{itemize}
\end{yourturn}

\noindent\textbf{(a)} Write a word problem for \quad \expr{-5 + 8 = 3}
\wlines{4}

\vspace{14pt}
\noindent\textbf{(b)} Write a word problem for \quad \expr{12 - (-4) = 16}
\wlines{4}

% ===========================================================================
% SECTION B -- SCIENCE
% ===========================================================================
\hwsection{Section B \textperiodcentered\ Science 1.1 --- Cells and Microscopes}

\begin{remember}
  A \key{cell} is the smallest basic unit of all living organisms. An \key{organelle}
  is a tiny structure inside a cell that does one specific, important job to keep the
  cell alive. Every cell has a \key{cell membrane}, \key{cytoplasm}, a \key{nucleus}
  and \key{mitochondria}. A plant cell has all of those \key{plus} a \key{cell wall}
  made of \key{cellulose}, \key{chloroplasts} containing \key{chlorophyll}, and a
  \key{sap vacuole}.
\end{remember}

% -- B1 ---------------------------------------------------------------------
\qhead{B1}{Match the word to its job}

\noindent Write the correct letter in the box next to each word.

\vspace{8pt}
\noindent
\begin{tabularx}{\textwidth}{@{}p{4.2cm} @{\hspace{6pt}} c @{\hspace{12pt}} X@{}}
  \textbf{Word} & \textbf{Letter} & \textbf{Job} \\[2pt]
  1.\ Cell            & \tickbox & \textbf{A} \; The control centre of the cell --- the ``boss''. \\
  2.\ Organelle       & \tickbox & \textbf{B} \; A large space of sugars and water that keeps a plant cell firm. \\
  3.\ Nucleus         & \tickbox & \textbf{C} \; Where energy is released from food. \\
  4.\ Mitochondria    & \tickbox & \textbf{D} \; The smallest basic unit of all living organisms. \\
  5.\ Cell wall       & \tickbox & \textbf{E} \; A tool that uses lenses to bend light and magnify an image. \\
  6.\ Chloroplast     & \tickbox & \textbf{F} \; A tiny structure inside a cell that does one specific job. \\
  7.\ Sap vacuole     & \tickbox & \textbf{G} \; A strong, stiff outer layer that holds a plant cell in shape. \\
  8.\ Microscope      & \tickbox & \textbf{H} \; Where a plant makes its food, using sunlight. \\
\end{tabularx}

% -- B2 ---------------------------------------------------------------------
\qhead{B2}{Animal, plant, or both?}

\noindent Put a tick ($\checkmark$) in \textbf{one} column for each organelle.

\vspace{8pt}
\noindent
\begin{tabularx}{\textwidth}{|X|>{\centering\arraybackslash}p{3.1cm}|>{\centering\arraybackslash}p{2.9cm}|>{\centering\arraybackslash}p{1.9cm}|}
  \hline
  \textbf{Organelle} & \textbf{Animal cells only} & \textbf{Plant cells only} & \textbf{Both} \\
  \hline
  Cell membrane & & & \\ \hline
  Cytoplasm     & & & \\ \hline
  Nucleus       & & & \\ \hline
  Mitochondria  & & & \\ \hline
  Cell wall     & & & \\ \hline
  Chloroplast   & & & \\ \hline
  Sap vacuole   & & & \\ \hline
\end{tabularx}

% -- B4 ---------------------------------------------------------------------
\qhead{B3}{Answer in full sentences}

\begin{wordhelp}[Sentence starters --- you may copy these]
  ``A plant cell needs a cell wall because\ldots''\\
  ``One difference is that\ldots\ Another difference is that\ldots''\\
  ``A microscope is a tool that\ldots''
\end{wordhelp}

\begin{enumerate}[label=\arabic*., itemsep=14pt]

  \item An animal has a skeleton to hold it up. A plant does not. Explain why every
        plant cell needs a cell wall.
        \wlines{2}

  \item Give \textbf{two} ways a cell wall is different from a cell membrane.
        Write each one as a full sentence.
        \wlines{4}

  \item Under the microscope, a leaf cell is full of small green circles. Name them,
        and name the green substance inside them.
        \wlines{2}

  \item What does a microscope do? Use the words \key{lens} and \key{magnify} in your
        answer.
        \wlines{3}

\end{enumerate}

% -- B5 ---------------------------------------------------------------------
\qhead{B4}{A model of a cell}

\noindent Mr Bowen builds a model of an animal cell in his kitchen. He uses a clear
plastic bag, filled with jelly, with a plum and four red beans pushed into the middle
of it.

\vspace{8pt}
\noindent\textbf{(a)} Which organelle does each item represent?

\vspace{6pt}
\noindent
\begin{tabularx}{\textwidth}{|p{5cm}|X|}
  \hline
  \textbf{Item in the model} & \textbf{Organelle it represents} \\
  \hline
  The plastic bag \rule{0pt}{1.9em} & \\ \hline
  The jelly       \rule{0pt}{1.9em} & \\ \hline
  The plum        \rule{0pt}{1.9em} & \\ \hline
  The four beans  \rule{0pt}{1.9em} & \\ \hline
\end{tabularx}

\vspace{14pt}
\noindent\textbf{(b)} Name \textbf{one} thing he could add to turn it into a model of a
\textbf{plant} cell, and say which organelle it would represent.
\wlines{2}

% -- B3 ---------------------------------------------------------------------
\qhead{B5}{Label the plant cell}

\noindent Write the name of each numbered part on the lines below the diagram.

\vspace{12pt}
\noindent
\resizebox{\textwidth}{!}{%
\begin{tikzpicture}[x=1cm,y=1cm,
    lbl/.style={circle, draw=cbteal, fill=white, line width=1.1pt,
                inner sep=1.8pt, minimum size=8mm, font=\bfseries}]

  % -- the cell itself. The wall band is deliberately 0.7 wide so that the
  %    wall (1) and the membrane (3) are visibly two different things.
  \fill[cbgreen!35, rounded corners=12pt] (0,0) rectangle (11,6.4);        % cell wall
  \fill[cbteal!10, rounded corners=8pt]  (0.7,0.7) rectangle (10.3,5.7);   % cytoplasm
  \draw[line width=1.2pt, black!60, rounded corners=8pt] (0.7,0.7) rectangle (10.3,5.7);
  \draw[line width=1.6pt, cbgreen!70!black, rounded corners=12pt] (0,0) rectangle (11,6.4);

  % -- sap vacuole --
  \fill[cbblue!18] (5.5,2.95) ellipse (3.2 and 1.65);
  \draw[line width=0.9pt, cbblue!60] (5.5,2.95) ellipse (3.2 and 1.65);

  % -- nucleus --
  \fill[cbpurple!30] (2.0,4.35) circle (0.85);
  \draw[line width=0.9pt, cbpurple!80] (2.0,4.35) circle (0.85);
  \fill[cbpurple!60] (2.0,4.35) circle (0.28);

  % -- chloroplasts --
  \foreach \x/\y/\r in {1.7/1.5/20, 3.2/1.15/-15, 5.0/1.05/10,
                        6.8/1.15/-20, 8.5/1.5/15, 9.6/2.6/-10, 9.6/4.2/25}{
    \begin{scope}[shift={(\x,\y)}, rotate=\r]
      \fill[cbgreen!75] (0,0) ellipse (0.46 and 0.24);
      \draw[line width=0.7pt, cbgreen!45!black] (0,0) ellipse (0.46 and 0.24);
    \end{scope}}

  % -- mitochondria. Kept away from x = 5.5 so the cytoplasm leader (2) has a
  %    clear, empty patch of cytoplasm to land on.
  \foreach \x/\y/\r in {3.3/5.1/-12, 8.4/5.1/16, 1.35/2.7/70}{
    \begin{scope}[shift={(\x,\y)}, rotate=\r]
      \fill[cborange!65] (0,0) ellipse (0.42 and 0.2);
      \draw[line width=0.7pt, cborange!80!black] (0,0) ellipse (0.42 and 0.2);
      \draw[line width=0.5pt, cborange!80!black] (-0.25,0.05) -- (-0.05,-0.06) -- (0.15,0.06);
    \end{scope}}

  % -- leader lines and numbers. Each one ends on a dot so there is no doubt
  %    about which structure it is pointing at.
  \foreach \sx/\sy/\ex/\ey/\lx/\ly/\n in {%
      2.2/6.4/2.2/7.55/2.2/7.95/1,          % cell wall  (the green band)
      5.5/5.35/7.9/7.55/8.0/7.95/2,         % cytoplasm  (empty patch)
      0.7/4.9/-1.5/6.25/-1.9/6.4/3,         % cell membrane (the black line)
      2.0/4.35/-1.5/4.05/-1.9/4.0/4,        % nucleus
      5.5/2.95/12.6/3.3/13.0/3.35/5,        % sap vacuole
      1.7/1.5/-1.5/1.1/-1.9/1.05/6,         % chloroplast
      8.4/5.1/12.6/5.5/13.0/5.55/7}{        % mitochondrion
    \draw[line width=0.7pt, black!65] (\sx,\sy) -- (\ex,\ey);
    \fill[black!65] (\sx,\sy) circle (0.075);
    \node[lbl] at (\lx,\ly) {\n};}

  % keep the bounding box sensible
  \path (-2.6,-0.5) rectangle (13.9,8.6);
\end{tikzpicture}}

\vspace{16pt}

% A tabular rather than multicols, so the seven answer lines can never be split
% across a page break away from their diagram.
\noindent
\begin{tabular}{@{}p{0.48\textwidth} p{0.48\textwidth}@{}}
  1.\ \blank[5.4cm] & 5.\ \blank[5.4cm] \\[30pt]
  2.\ \blank[5.4cm] & 6.\ \blank[5.4cm] \\[30pt]
  3.\ \blank[5.4cm] & 7.\ \blank[5.4cm] \\[30pt]
  4.\ \blank[5.4cm] &                   \\
\end{tabular}

\vspace{10pt}
\noindent{\small\color{black!60}Hint: number 3 is the very thin layer just inside the
outer green edge. Numbers 1 and 3 are easy to confuse --- one is thick and stiff, one is
thin and flexible.}

% ===========================================================================
% ANSWER KEY (teacher copy only)
% ===========================================================================
\ifkey
\newpage
\hwsection[cbred]{Answer Key --- Teacher Copy}

% The key is for an adult reading it at a desk, not a 12-year-old writing on it,
% so tighten everything back up.
\linespread{1.0}\selectfont
\setlist[enumerate]{itemsep=2pt, topsep=3pt, leftmargin=*}
\setlist[itemize]{itemsep=2pt, topsep=3pt, leftmargin=*}

\qhead[cbred]{A1}{Show it on the number line}
\noindent (a) Cross on $4$, arrow three steps \emph{left}, landing on $1$. So
$4+(-3)=1$. The cross must be on 4, not on 3 --- adding a negative means starting at
the first number and moving left.\\
(b) (i) lower than \quad (ii) warmer than \quad (iii) lower than.
{\small\ (``higher than'' is also acceptable in (iii) if the sentence is reversed ---
mark the reasoning, not the word.)}

\qhead[cbred]{A2}{Moving along the line}
\noindent
(a) $-2$ \quad (b) $-4$ \quad (c) $-8$ \quad (d) $5$ \quad (e) $0$ \quad (f) $-6$ \\
(g) $6$ \quad (h) $8$ \quad (i) $-7$ \quad (j) $0$ \quad (k) $3$ \quad (l) $-16$

\smallskip
\noindent{\small Common errors: (d) and (f) --- students flip \emph{any} pair of signs
to $+$. Only $-(-)$ becomes $+$. (b) $5+(-9)$ is a single $+(-)$, so it still goes
left.}

\qhead[cbred]{A3}{Say it, then write it}
\begin{enumerate}[label=\arabic*., itemsep=3pt]
  \item $4-9=-5$
  \item $-2-6=-8$
  \item $3-7=-4$
  \item $-8+5=-3$
  \item $-30-12=-42$
  \item $-35-(-20)=-15$
\end{enumerate}
\noindent{\small Questions 1--3 are the ``reversed order'' trap; 6 is the ``removing a
debt adds'' rule --- the only place it is tested now, so it is worth checking. If a student writes $9-4$, $6-(-2)$ or $7-3$, that is the English
error, not an arithmetic one --- say so, and make them re-read the sentence.}

\qhead[cbred]{A4}{Difference}
\begin{enumerate}[label=\arabic*., itemsep=3pt]
  \item $8$
  \item $6$
  \item $4-(-13)=17$ degrees
  \item (a) $11$ \; (b) $-11$ \; (c) A difference is the gap between the numbers so it
        is never negative, but $-5-6$ is an instruction to move 6 further left.
  \item $18-(-24)=42$ degrees
\end{enumerate}

\qhead[cbred]{A5}{Word problems}
\begin{enumerate}[label=\arabic*., itemsep=3pt]
  \item $7-11+2=-2$. Floor $-2$ (B2).
  \item (a) $-3+11-14=-6\dC$. \; (b) Warmest $8\dC$ at noon, coldest $-6\dC$:
        $8-(-6)=14$ degrees.
  \item Still $-40$ dollars. Each week $+25-25=0$, so six weeks change nothing. This one
        rewards reading to the end before calculating.
\end{enumerate}

\qhead[cbred]{A6}{Now you write the question}
\noindent There is no single right answer --- check the four conditions in the green
box, in this order:
\begin{enumerate}[label=\arabic*., itemsep=3pt]
  \item \textbf{Does it start in the right place?} For (a) the story must \emph{begin}
        at $-5$, not arrive there. This is the same ``which number comes first''
        problem as A3, and it is the one most of them will get wrong.
  \item \textbf{Does the signal word match the sign?} (a) needs an \emph{up} word
        (rises, deposits, climbs) for the $+8$. (b) needs a \emph{removing a negative}
        --- a debt cancelled, a weight taken off --- not simply ``adds 4''.
  \item \textbf{Full sentences, ending in a question?}
  \item \textbf{Does their own answer come out as 3, and as 16?}
\end{enumerate}

\smallskip
\noindent{\small Two that work: (a) ``The temperature in Sa Pa was $-5\dC$ at
midnight. By noon it had risen by 8 degrees. What was the temperature at noon?''
(b) ``Mr Bowen had 12 dollars, and he also owed his brother 4 dollars. His brother told
him to forget the debt. How much is Mr Bowen worth now?'' \\
Part (b) is the harder one by a distance --- expect most of the class to write a plain
addition. That is worth a whole-class discussion rather than a red pen: ask them what
in real life makes you \emph{richer} when it is taken away.}

\qhead[cbred]{B1}{Match the word to its job}
\noindent 1--D \quad 2--F \quad 3--A \quad 4--C \quad 5--G \quad 6--H \quad 7--B \quad 8--E

\qhead[cbred]{B2}{Animal, plant, or both?}
\noindent Both: cell membrane, cytoplasm, nucleus, mitochondria. \\
Plant only: cell wall, chloroplast, sap vacuole. \\
{\small Nothing is animal-only. Expect some students to tick that column anyway ---
it is worth saying out loud that a plant cell has everything an animal cell has,
\emph{plus} extras.}

\qhead[cbred]{B3}{Answer in full sentences}
\begin{enumerate}[label=\arabic*., itemsep=3pt]
  \item A plant has no skeleton, so each cell needs a stiff wall around it to hold its
        shape and hold the plant up.
  \item Any two: the wall is thick and stiff while the membrane is thin and flexible;
        the wall is made of cellulose; the wall is only in plant cells while every cell
        has a membrane; the membrane controls what enters and leaves, the wall does not.
  \item They are chloroplasts, and the green substance inside them is chlorophyll.
  \item A microscope is a tool that uses curved pieces of glass called lenses to bend
        light and magnify an image, making it look much bigger than it really is.
\end{enumerate}

\qhead[cbred]{B4}{A model of a cell}
\noindent (a) bag --- cell membrane; jelly --- cytoplasm; plum --- nucleus;
beans --- mitochondria.

\smallskip
\noindent (b) Any one, with its organelle named: a shoebox or rigid container around the
bag --- cell wall; a water balloon in the middle --- sap vacuole; green sweets ---
chloroplasts.

\qhead[cbred]{B5}{Label the plant cell}
\noindent 1 --- cell wall \quad 2 --- cytoplasm \quad 3 --- cell membrane \quad
4 --- nucleus \quad 5 --- sap vacuole \quad 6 --- chloroplast \quad
7 --- mitochondrion

\fi

\end{document}
"

\documentclass[12pt,a4paper]{article}

% -- Answer key switch -------------------------------------------------------
\newif\ifkey
\ifdefined\TEACHER\keytrue\else
  \keyfalse        % \keyfalse = student copy   |   \keytrue = teacher copy
\fi

% -- House style -------------------------------------------------------------
% Shared by every packet. See homework/hw-style.tex.
% homework/hw-style.tex
% Shared house style for every Year 7 homework packet. \input this from the
% packet file, right after \documentclass and the answer-key switch.
%
% Do not fork this file per packet. If a packet needs a different look, the
% question is whether every packet should have it.
%
% The choices here are deliberate and were arrived at by printing and looking:
%   * No solid dark banners. Every header is a light tint with a coloured spine
%     on the left, so colour reads clearly without flooding a school printer.
%   * No marks and no mark totals. These are practice, not exam papers.
%   * Lato at 12pt with open leading, plus mathastext so inline maths is Lato
%     too. Without mathastext every "-6 + 4" is Computer Modern serif and the
%     sheet looks half-typeset.
%   * Writing lines are spaced for Year 7 handwriting, not adult margin notes.
%
% Provides: \wlines \blank \tickbox \sep \qmk-free marks (none) \hwsection
% \qhead \expr, the workspace box, and the remember / wordhelp / activity /
% yourturn coloured boxes. Palette matches the lesson decks exactly.

% -- Packages ----------------------------------------------------------------
\usepackage[T1]{fontenc}
\usepackage[utf8]{inputenc}
\usepackage[default]{lato}
\usepackage[a4paper,top=1.7cm,bottom=1.7cm,left=2.0cm,right=2.0cm,headheight=15pt]{geometry}
\usepackage{amsmath,amssymb}
\usepackage{xcolor}
\usepackage[most]{tcolorbox}
\usepackage{enumitem}
\usepackage{tikz}
\usetikzlibrary{arrows.meta}
\usepackage{array}
\usepackage{tabularx}
\usepackage{fancyhdr}
\usepackage{multicol}
\usepackage{needspace}
\usepackage{graphicx}

% Make maths use Lato too. Without this, every "-6 + 4" on the page is set in
% Computer Modern serif and the sheet looks half-typeset. Must come last.
\usepackage[italic]{mathastext}

\linespread{1.06}\selectfont

% -- The Learner's Book palette (same hex values as the lesson decks) ---------
\definecolor{cbteal}{HTML}{0087A8}
\definecolor{cbpurple}{HTML}{5C2483}
\definecolor{cborange}{HTML}{C25E12}
\definecolor{cbgreen}{HTML}{4A8B23}
\definecolor{cbred}{HTML}{C8102E}
\definecolor{cbblue}{HTML}{1A5FA8}
\definecolor{rulegray}{HTML}{9AA5AD}

% -- Answer lines and blanks -------------------------------------------------
% \wlines{n} draws n ruled writing lines, spaced wide enough for Year 7
% handwriting rather than for an adult with a fine pen.
\makeatletter
\newcommand{\wlines}[1]{%
  \par\vspace{0.75em}%
  \@tempcnta=0\relax
  \@whilenum\@tempcnta<#1\do{%
    \hbox to \linewidth{\color{rulegray}\leaders\hrule height 0.5pt\hfill}%
    \vspace{1.5em}%
    \advance\@tempcnta by 1\relax}%
  \par\vspace{-0.8em}%
}
\makeatother

% \blank[width] is an inline gap to write a short answer in.
\newcommand{\blank}[1][2.6cm]{\underline{\hspace{#1}}}

% A boxed space for showing working.
\newtcolorbox{workspace}[1][3.4cm]{%
  colback=white, colframe=rulegray, boxrule=0.5pt, arc=5pt,
  height=#1, left=6pt, right=6pt, top=4pt, bottom=4pt,
  before skip=9pt, after skip=11pt}

% An empty tick box.
\newcommand{\tickbox}{\fbox{\rule{0pt}{1.15em}\rule{1.15em}{0pt}}}

% A middot separator.
\newcommand{\sep}{\,\textperiodcentered\,}

% -- Section and question headings -------------------------------------------
% Light tint plus a fat coloured spine on the left. No solid dark banner: the
% colour reads clearly but barely costs the printer anything.
% needspace stops a heading stranding alone at the foot of a page. Kept modest:
% too large and it pushes headings over, leaving a worse hole than it prevents.
\newcommand{\hwsection}[2][cbteal]{%
  \needspace{8\baselineskip}%
  \par\vspace{13pt}%
  \noindent
  \begin{tcolorbox}[enhanced, nobeforeafter, width=\textwidth,
      colback=#1!7, colframe=#1!7, arc=5pt,
      borderline west={5pt}{0pt}{#1},
      left=14pt, right=10pt, top=7pt, bottom=7pt]
    {\large\bfseries\color{#1}#2}
  \end{tcolorbox}%
  \par\vspace{9pt}}

% \qhead[colour]{A2}{Moving along the line} -- a question heading.
\newcommand{\qhead}[3][cbteal]{%
  \needspace{6\baselineskip}%
  \par\vspace{11pt}%
  \noindent{\large\bfseries\color{#1}#2}\quad{\large\bfseries #3}\par\vspace{4pt}}

% -- Coloured boxes ----------------------------------------------------------
% Shared look: rounded, light fill, coloured rule, coloured title text on a
% pale title strip rather than a dark one.
\tcbset{hwbox/.style={
  enhanced, breakable, arc=6pt, boxrule=1pt,
  left=11pt, right=11pt, top=8pt, bottom=8pt,
  fonttitle=\bfseries, attach boxed title to top left={xshift=12pt, yshift=-8pt},
  boxed title style={arc=4pt, boxrule=1pt, left=7pt, right=7pt, top=2pt, bottom=2pt},
  top=13pt}}

% Orange: things to remember. Matches the deck's "Write This Down".
\newtcolorbox{remember}[1][Remember from the lesson]{hwbox,
  colback=cborange!6, colframe=cborange, coltitle=cborange,
  colbacktitle=white, title={#1}}

% Blue: English help. The barrier in this class is the wording, not the maths.
\newtcolorbox{wordhelp}[1][Word help]{hwbox,
  colback=cbblue!6, colframe=cbblue, coltitle=cbblue,
  colbacktitle=white, title={#1}}

% Purple: an activity or instruction.
\newtcolorbox{activity}[1][How to do this homework]{hwbox,
  colback=cbpurple!6, colframe=cbpurple, coltitle=cbpurple,
  colbacktitle=white, title={#1}}

% Green: your turn / a friendly nudge.
\newtcolorbox{yourturn}[1][Your turn]{hwbox,
  colback=cbgreen!7, colframe=cbgreen, coltitle=cbgreen!75!black,
  colbacktitle=white, title={#1}}

% -- Shorthands --------------------------------------------------------------
\newcommand{\dC}{\,^{\circ}\mathrm{C}}          % use inside maths: $-8\dC$
\newcommand{\key}[1]{{\color{cborange}\textbf{#1}}}
\newcommand{\eng}[1]{{\color{cbblue}\textbf{#1}}}

\setlist[enumerate]{itemsep=8pt, topsep=5pt, leftmargin=*}
\setlist[itemize]{itemsep=4pt, topsep=4pt, leftmargin=*}
\renewcommand{\arraystretch}{1.45}

% -- An expression the student is asked to work from -------------------------
\newcommand{\expr}[1]{%
  \tcbox[on line, colback=cbgreen!10, colframe=cbgreen!55, boxrule=1pt, arc=4pt,
         left=9pt, right=9pt, top=4pt, bottom=4pt]{\large\bfseries $#1$}}

% -- Running head and foot ---------------------------------------------------
% The packet overrides \fancyhead[L] with its own title.
\pagestyle{fancy}
\fancyhf{}
\fancyhead[L]{\small\color{black!55}Year 7\sep Homework}
\fancyhead[R]{\small\color{black!55}Name: \blank[3.4cm]}
\fancyfoot[C]{\small\color{black!55}Page \thepage}
\renewcommand{\headrulewidth}{0pt}
\renewcommand{\headrule}{\vspace{2pt}\hbox to\headwidth{\color{cbteal!45}\leaders\hrule height 1.2pt\hfill}}



% -- Per-packet running head -------------------------------------------------
\fancyhead[L]{\small\color{black!55}Year 7\sep Homework 1}

\begin{document}

% ===========================================================================
% COVER BLOCK
% ===========================================================================
\thispagestyle{fancy}

\begin{tcolorbox}[enhanced, colback=cbpurple!7, colframe=cbpurple!7, arc=7pt,
                  borderline west={6pt}{0pt}{cbpurple},
                  left=18pt, right=14pt, top=14pt, bottom=14pt]
  {\color{cbpurple}\bfseries YEAR 7\sep UNIT 1\sep HOMEWORK 1}\\[6pt]
  {\color{cbpurple}\Huge\bfseries Integers and Cells}\\[10pt]
  {\color{black!70}Mathematics 1.1 --- Adding \& Subtracting Integers}\\
  {\color{black!70}Science 1.1 --- Cells and Microscopes}
\end{tcolorbox}

\vspace{10pt}
\noindent
\begin{tabularx}{\textwidth}{@{}X X@{}}
  \textbf{Name:} \blank[5.6cm] & \textbf{Class:} \blank[5.6cm] \\
  \textbf{Date given:} \blank[4.6cm] & \textbf{Date due:} \blank[5.0cm] \\
\end{tabularx}

\vspace{8pt}
\begin{activity}
  \begin{itemize}
    \item Write your answers \textbf{on this sheet}, in the spaces given. If you need
          more room, use the back of the page.
    \item In maths, \key{write the calculation before the answer}.
    \item In science, answer in \key{full English sentences}. ``Chloroplast'' is not a
          sentence; ``The green circles are chloroplasts'' is.
    \item If a word is new to you, look for the blue \eng{Word help} box on that page
          before you give up.
  \end{itemize}
\end{activity}

% ===========================================================================
% SECTION A -- MATHEMATICS
% ===========================================================================
\hwsection{Section A \textperiodcentered\ Mathematics 1.1 --- Adding and Subtracting Integers}

\begin{remember}
  An \key{integer} is a whole number that is positive, negative or zero --- never a
  fraction. On the number line, \key{negative} integers sit to the \key{left} of zero
  and \key{positive} integers to the \key{right}.

  \medskip
  \begin{tabularx}{\linewidth}{@{}lX@{}}
    Add a \key{positive}      & move \key{right} \\
    Add a \key{negative}      & move \key{left} \\
    Subtract a \key{positive} & move \key{left} \\
    Subtract a \key{negative} & move \key{right} --- because $a-(-b)=a+b$ \\
  \end{tabularx}

  \medskip
  Careful: only $-(-)$ becomes $+$. A single $+(-)$ still sends you left, so
  $5+(-3)=2$, \textbf{not} $8$.
\end{remember}

% -- A1 ---------------------------------------------------------------------
\qhead{A1}{Show it on the number line}

\noindent\textbf{(a)} Show $4 + (-3)$ on the number line below. Put a cross on the
number you \textbf{start} at, then draw an arrow for the move.

\vspace{12pt}
\noindent
\resizebox{0.92\textwidth}{!}{%
\begin{tikzpicture}[x=1cm,y=1cm]
  \draw[line width=1.1pt,cbteal,{Latex[length=3mm]}-{Latex[length=3mm]}] (-5.7,0) -- (8.7,0);
  \foreach \n in {-5,...,8}{
    \draw[line width=0.9pt,cbteal] (\n,-0.18) -- (\n,0.18);
    \node[below=4pt,font=\small] at (\n,-0.18) {$\n$};
  }
\end{tikzpicture}}

\vspace{10pt}
\noindent\hspace*{0.5em}$4 + (-3) = \blank[2.2cm]$

\vspace{12pt}
\noindent\textbf{(b)} Complete each sentence with \eng{higher than}, \eng{lower than},
\eng{warmer than} or \eng{colder than}.

\begin{enumerate}[label=(\roman*)]
  \item $-8$ is \blank[3.4cm] $-4$.
  \item A temperature of $3\dC$ is \blank[3.4cm] a temperature of $-9\dC$.
  \item Floor $-2$ of the car park is \blank[3.4cm] floor $1$.
\end{enumerate}

% -- A2 ---------------------------------------------------------------------
\qhead{A2}{Moving along the line}

\noindent Work these out. You may draw a small number line beside any of them.

\vspace{10pt}
\begin{multicols}{3}
\begin{enumerate}[label=(\alph*), itemsep=13pt]
  \item $-6+4 = \blank[1.6cm]$
  \item $5+(-9) = \blank[1.6cm]$
  \item $-3-5 = \blank[1.6cm]$
  \item $-2-(-7) = \blank[1.6cm]$
  \item $8+(-8) = \blank[1.6cm]$
  \item $-10-(-4) = \blank[1.6cm]$
  \item $0-(-6) = \blank[1.6cm]$
  \item $-7+15 = \blank[1.6cm]$
  \item $4-11 = \blank[1.6cm]$
  \item $-5-(-5) = \blank[1.6cm]$
  \item $6-8-(-5) = \blank[1.6cm]$
  \item $-9-6+(-1) = \blank[1.6cm]$
\end{enumerate}
\end{multicols}

% -- A3 ---------------------------------------------------------------------
\qhead{A3}{Say it, then write it}

\begin{wordhelp}[Word help --- these words decide the order of your calculation]
  \begin{tabularx}{\linewidth}{@{}lX@{}}
    \eng{subtract 5 from 8}   & you \textbf{start} at 8: $8-5$. The English gives the numbers in the \textbf{opposite order} to the calculation. \\
    \eng{take 6 away from 2}  & the same trap: $2-6$. \\
    \eng{7 less than 3}       & again the same trap: $3-7$. \\
    \eng{rises by / falls by} & the temperature goes up / the temperature goes down. \\
    \eng{deposit}             & to put money \textbf{into} an account $(+)$. \\
    \eng{withdraw}            & to take money \textbf{out of} an account $(-)$. \\
    \eng{owe} / \eng{a debt}  & money you must pay back. It makes your total \textbf{negative}. \\
    \eng{descend}             & to go \textbf{down}. To go up is to \eng{ascend}. \\
  \end{tabularx}
\end{wordhelp}

\noindent Turn each English sentence into a calculation, then work out the answer.
\key{Write the calculation first.}

\vspace{6pt}
\begin{enumerate}[label=\arabic*., itemsep=11pt]
  \item Subtract 9 from 4. \\[4pt]
        Calculation: \blank[6.2cm] \quad Answer: \blank[2.4cm]
  \item Take 6 away from $-2$. \\[4pt]
        Calculation: \blank[6.2cm] \quad Answer: \blank[2.4cm]
  \item What is 7 less than 3? \\[4pt]
        Calculation: \blank[6.2cm] \quad Answer: \blank[2.4cm]
  \item The temperature is $-8\dC$ and it rises by 5 degrees. \\[4pt]
        Calculation: \blank[6.2cm] \quad Answer: \blank[2.4cm]
  \item A submarine is at $-30$ m. It descends a further 12 m. \\[4pt]
        Calculation: \blank[6.2cm] \quad Answer: \blank[2.4cm]
  \item An account holds $-35$ dollars. The bank removes a debt of 20 dollars from it.
        \\[4pt]
        Calculation: \blank[6.2cm] \quad Answer: \blank[2.4cm]
\end{enumerate}

% -- A4 ---------------------------------------------------------------------
\qhead{A4}{Difference}

\begin{remember}[Remember: a difference is a gap]
  The \key{difference} between two numbers is how far apart they are on the line.
  Work it out with \key{bigger $-$ smaller}. A difference is \key{never negative}.
  \eng{How much warmer}, \eng{how much colder}, \eng{how much higher} and
  \eng{how many more} all ask for the same thing.
\end{remember}

\begin{enumerate}[label=\arabic*.]
  \item Find the difference between $-6$ and $2$. \hfill Answer: \blank[2.8cm]
  \item Find the difference between $-9$ and $-3$. \hfill Answer: \blank[2.8cm]
  \item How much warmer is $4\dC$ than $-13\dC$? Show your calculation.
        \wlines{1}
  \item These two questions look almost the same. They are not.
        \begin{enumerate}[label=(\alph*), itemsep=9pt]
          \item Find the difference between $-5$ and $6$. \hfill Answer: \blank[2.8cm]
          \item What is $-5$ minus $6$? \hfill Answer: \blank[2.8cm]
          \item In one full sentence, explain why the two answers are different.
                \wlines{2}
        \end{enumerate}
  \item One winter morning it is $18\dC$ in Ho Chi Minh City and $-24\dC$ in
        Ulaanbaatar. How much colder is Ulaanbaatar? Show your calculation.
        \wlines{1}
\end{enumerate}

% -- A5 ---------------------------------------------------------------------
\qhead{A5}{Word problems}

\noindent Read each one \key{all the way to the end} before you start writing.
Show the calculation, then the answer.

\begin{enumerate}[label=\textbf{\arabic*.}, itemsep=16pt]

  \item \textbf{The office lift.} Mr Bowen gets into the lift on floor 7. He goes
        \textbf{down} 11 floors, and then \textbf{up} 2 floors. Which floor is he on
        now?
        \begin{workspace}[2.8cm]\end{workspace}

  \item \textbf{A day in Sa Pa.} At 6 a.m. the temperature is $-3\dC$. By noon it has
        risen 11 degrees. By midnight it has fallen 14 degrees.
        \begin{enumerate}[label=(\alph*)]
          \item What is the temperature at midnight?
          \item What is the \textbf{difference} between the warmest and the coldest
                temperature of the day?
        \end{enumerate}
        \begin{workspace}[3.4cm]\end{workspace}

  \item \textbf{Six weeks of saving.} Mr Bowen's account is at $-40$ dollars. Every
        Monday he deposits 25 dollars. Every Friday he withdraws 25 dollars. How much is
        in the account after 6 weeks? Explain your answer in one full sentence.
        \begin{workspace}[3.1cm]\end{workspace}

\end{enumerate}

% -- A6 ---------------------------------------------------------------------
\qhead[cbgreen]{A6}{Now you write the question}

\begin{yourturn}[Your turn to be the teacher]
  So far the questions have given you the English, and you have written the maths.
  Now do it \textbf{backwards}. You are given the maths --- write the English.

  \medskip
  Your word problem must:
  \begin{itemize}
    \item be about something \textbf{real} --- money, temperature, floors in a
          building, or depth under the sea;
    \item use at least one word from the \eng{Word help} box on the earlier page,
          such as \eng{rises by}, \eng{falls by}, \eng{deposits}, \eng{withdraws},
          \eng{owes}, \eng{below} or \eng{descends};
    \item be written in \textbf{full English sentences}, and end with a question;
    \item have the answer that is already given to you.
  \end{itemize}
\end{yourturn}

\noindent\textbf{(a)} Write a word problem for \quad \expr{-5 + 8 = 3}
\wlines{4}

\vspace{14pt}
\noindent\textbf{(b)} Write a word problem for \quad \expr{12 - (-4) = 16}
\wlines{4}

% ===========================================================================
% SECTION B -- SCIENCE
% ===========================================================================
\hwsection{Section B \textperiodcentered\ Science 1.1 --- Cells and Microscopes}

\begin{remember}
  A \key{cell} is the smallest basic unit of all living organisms. An \key{organelle}
  is a tiny structure inside a cell that does one specific, important job to keep the
  cell alive. Every cell has a \key{cell membrane}, \key{cytoplasm}, a \key{nucleus}
  and \key{mitochondria}. A plant cell has all of those \key{plus} a \key{cell wall}
  made of \key{cellulose}, \key{chloroplasts} containing \key{chlorophyll}, and a
  \key{sap vacuole}.
\end{remember}

% -- B1 ---------------------------------------------------------------------
\qhead{B1}{Match the word to its job}

\noindent Write the correct letter in the box next to each word.

\vspace{8pt}
\noindent
\begin{tabularx}{\textwidth}{@{}p{4.2cm} @{\hspace{6pt}} c @{\hspace{12pt}} X@{}}
  \textbf{Word} & \textbf{Letter} & \textbf{Job} \\[2pt]
  1.\ Cell            & \tickbox & \textbf{A} \; The control centre of the cell --- the ``boss''. \\
  2.\ Organelle       & \tickbox & \textbf{B} \; A large space of sugars and water that keeps a plant cell firm. \\
  3.\ Nucleus         & \tickbox & \textbf{C} \; Where energy is released from food. \\
  4.\ Mitochondria    & \tickbox & \textbf{D} \; The smallest basic unit of all living organisms. \\
  5.\ Cell wall       & \tickbox & \textbf{E} \; A tool that uses lenses to bend light and magnify an image. \\
  6.\ Chloroplast     & \tickbox & \textbf{F} \; A tiny structure inside a cell that does one specific job. \\
  7.\ Sap vacuole     & \tickbox & \textbf{G} \; A strong, stiff outer layer that holds a plant cell in shape. \\
  8.\ Microscope      & \tickbox & \textbf{H} \; Where a plant makes its food, using sunlight. \\
\end{tabularx}

% -- B2 ---------------------------------------------------------------------
\qhead{B2}{Animal, plant, or both?}

\noindent Put a tick ($\checkmark$) in \textbf{one} column for each organelle.

\vspace{8pt}
\noindent
\begin{tabularx}{\textwidth}{|X|>{\centering\arraybackslash}p{3.1cm}|>{\centering\arraybackslash}p{2.9cm}|>{\centering\arraybackslash}p{1.9cm}|}
  \hline
  \textbf{Organelle} & \textbf{Animal cells only} & \textbf{Plant cells only} & \textbf{Both} \\
  \hline
  Cell membrane & & & \\ \hline
  Cytoplasm     & & & \\ \hline
  Nucleus       & & & \\ \hline
  Mitochondria  & & & \\ \hline
  Cell wall     & & & \\ \hline
  Chloroplast   & & & \\ \hline
  Sap vacuole   & & & \\ \hline
\end{tabularx}

% -- B4 ---------------------------------------------------------------------
\qhead{B3}{Answer in full sentences}

\begin{wordhelp}[Sentence starters --- you may copy these]
  ``A plant cell needs a cell wall because\ldots''\\
  ``One difference is that\ldots\ Another difference is that\ldots''\\
  ``A microscope is a tool that\ldots''
\end{wordhelp}

\begin{enumerate}[label=\arabic*., itemsep=14pt]

  \item An animal has a skeleton to hold it up. A plant does not. Explain why every
        plant cell needs a cell wall.
        \wlines{2}

  \item Give \textbf{two} ways a cell wall is different from a cell membrane.
        Write each one as a full sentence.
        \wlines{4}

  \item Under the microscope, a leaf cell is full of small green circles. Name them,
        and name the green substance inside them.
        \wlines{2}

  \item What does a microscope do? Use the words \key{lens} and \key{magnify} in your
        answer.
        \wlines{3}

\end{enumerate}

% -- B5 ---------------------------------------------------------------------
\qhead{B4}{A model of a cell}

\noindent Mr Bowen builds a model of an animal cell in his kitchen. He uses a clear
plastic bag, filled with jelly, with a plum and four red beans pushed into the middle
of it.

\vspace{8pt}
\noindent\textbf{(a)} Which organelle does each item represent?

\vspace{6pt}
\noindent
\begin{tabularx}{\textwidth}{|p{5cm}|X|}
  \hline
  \textbf{Item in the model} & \textbf{Organelle it represents} \\
  \hline
  The plastic bag \rule{0pt}{1.9em} & \\ \hline
  The jelly       \rule{0pt}{1.9em} & \\ \hline
  The plum        \rule{0pt}{1.9em} & \\ \hline
  The four beans  \rule{0pt}{1.9em} & \\ \hline
\end{tabularx}

\vspace{14pt}
\noindent\textbf{(b)} Name \textbf{one} thing he could add to turn it into a model of a
\textbf{plant} cell, and say which organelle it would represent.
\wlines{2}

% -- B3 ---------------------------------------------------------------------
\qhead{B5}{Label the plant cell}

\noindent Write the name of each numbered part on the lines below the diagram.

\vspace{12pt}
\noindent
\resizebox{\textwidth}{!}{%
\begin{tikzpicture}[x=1cm,y=1cm,
    lbl/.style={circle, draw=cbteal, fill=white, line width=1.1pt,
                inner sep=1.8pt, minimum size=8mm, font=\bfseries}]

  % -- the cell itself. The wall band is deliberately 0.7 wide so that the
  %    wall (1) and the membrane (3) are visibly two different things.
  \fill[cbgreen!35, rounded corners=12pt] (0,0) rectangle (11,6.4);        % cell wall
  \fill[cbteal!10, rounded corners=8pt]  (0.7,0.7) rectangle (10.3,5.7);   % cytoplasm
  \draw[line width=1.2pt, black!60, rounded corners=8pt] (0.7,0.7) rectangle (10.3,5.7);
  \draw[line width=1.6pt, cbgreen!70!black, rounded corners=12pt] (0,0) rectangle (11,6.4);

  % -- sap vacuole --
  \fill[cbblue!18] (5.5,2.95) ellipse (3.2 and 1.65);
  \draw[line width=0.9pt, cbblue!60] (5.5,2.95) ellipse (3.2 and 1.65);

  % -- nucleus --
  \fill[cbpurple!30] (2.0,4.35) circle (0.85);
  \draw[line width=0.9pt, cbpurple!80] (2.0,4.35) circle (0.85);
  \fill[cbpurple!60] (2.0,4.35) circle (0.28);

  % -- chloroplasts --
  \foreach \x/\y/\r in {1.7/1.5/20, 3.2/1.15/-15, 5.0/1.05/10,
                        6.8/1.15/-20, 8.5/1.5/15, 9.6/2.6/-10, 9.6/4.2/25}{
    \begin{scope}[shift={(\x,\y)}, rotate=\r]
      \fill[cbgreen!75] (0,0) ellipse (0.46 and 0.24);
      \draw[line width=0.7pt, cbgreen!45!black] (0,0) ellipse (0.46 and 0.24);
    \end{scope}}

  % -- mitochondria. Kept away from x = 5.5 so the cytoplasm leader (2) has a
  %    clear, empty patch of cytoplasm to land on.
  \foreach \x/\y/\r in {3.3/5.1/-12, 8.4/5.1/16, 1.35/2.7/70}{
    \begin{scope}[shift={(\x,\y)}, rotate=\r]
      \fill[cborange!65] (0,0) ellipse (0.42 and 0.2);
      \draw[line width=0.7pt, cborange!80!black] (0,0) ellipse (0.42 and 0.2);
      \draw[line width=0.5pt, cborange!80!black] (-0.25,0.05) -- (-0.05,-0.06) -- (0.15,0.06);
    \end{scope}}

  % -- leader lines and numbers. Each one ends on a dot so there is no doubt
  %    about which structure it is pointing at.
  \foreach \sx/\sy/\ex/\ey/\lx/\ly/\n in {%
      2.2/6.4/2.2/7.55/2.2/7.95/1,          % cell wall  (the green band)
      5.5/5.35/7.9/7.55/8.0/7.95/2,         % cytoplasm  (empty patch)
      0.7/4.9/-1.5/6.25/-1.9/6.4/3,         % cell membrane (the black line)
      2.0/4.35/-1.5/4.05/-1.9/4.0/4,        % nucleus
      5.5/2.95/12.6/3.3/13.0/3.35/5,        % sap vacuole
      1.7/1.5/-1.5/1.1/-1.9/1.05/6,         % chloroplast
      8.4/5.1/12.6/5.5/13.0/5.55/7}{        % mitochondrion
    \draw[line width=0.7pt, black!65] (\sx,\sy) -- (\ex,\ey);
    \fill[black!65] (\sx,\sy) circle (0.075);
    \node[lbl] at (\lx,\ly) {\n};}

  % keep the bounding box sensible
  \path (-2.6,-0.5) rectangle (13.9,8.6);
\end{tikzpicture}}

\vspace{16pt}

% A tabular rather than multicols, so the seven answer lines can never be split
% across a page break away from their diagram.
\noindent
\begin{tabular}{@{}p{0.48\textwidth} p{0.48\textwidth}@{}}
  1.\ \blank[5.4cm] & 5.\ \blank[5.4cm] \\[30pt]
  2.\ \blank[5.4cm] & 6.\ \blank[5.4cm] \\[30pt]
  3.\ \blank[5.4cm] & 7.\ \blank[5.4cm] \\[30pt]
  4.\ \blank[5.4cm] &                   \\
\end{tabular}

\vspace{10pt}
\noindent{\small\color{black!60}Hint: number 3 is the very thin layer just inside the
outer green edge. Numbers 1 and 3 are easy to confuse --- one is thick and stiff, one is
thin and flexible.}

% ===========================================================================
% ANSWER KEY (teacher copy only)
% ===========================================================================
\ifkey
\newpage
\hwsection[cbred]{Answer Key --- Teacher Copy}

% The key is for an adult reading it at a desk, not a 12-year-old writing on it,
% so tighten everything back up.
\linespread{1.0}\selectfont
\setlist[enumerate]{itemsep=2pt, topsep=3pt, leftmargin=*}
\setlist[itemize]{itemsep=2pt, topsep=3pt, leftmargin=*}

\qhead[cbred]{A1}{Show it on the number line}
\noindent (a) Cross on $4$, arrow three steps \emph{left}, landing on $1$. So
$4+(-3)=1$. The cross must be on 4, not on 3 --- adding a negative means starting at
the first number and moving left.\\
(b) (i) lower than \quad (ii) warmer than \quad (iii) lower than.
{\small\ (``higher than'' is also acceptable in (iii) if the sentence is reversed ---
mark the reasoning, not the word.)}

\qhead[cbred]{A2}{Moving along the line}
\noindent
(a) $-2$ \quad (b) $-4$ \quad (c) $-8$ \quad (d) $5$ \quad (e) $0$ \quad (f) $-6$ \\
(g) $6$ \quad (h) $8$ \quad (i) $-7$ \quad (j) $0$ \quad (k) $3$ \quad (l) $-16$

\smallskip
\noindent{\small Common errors: (d) and (f) --- students flip \emph{any} pair of signs
to $+$. Only $-(-)$ becomes $+$. (b) $5+(-9)$ is a single $+(-)$, so it still goes
left.}

\qhead[cbred]{A3}{Say it, then write it}
\begin{enumerate}[label=\arabic*., itemsep=3pt]
  \item $4-9=-5$
  \item $-2-6=-8$
  \item $3-7=-4$
  \item $-8+5=-3$
  \item $-30-12=-42$
  \item $-35-(-20)=-15$
\end{enumerate}
\noindent{\small Questions 1--3 are the ``reversed order'' trap; 6 is the ``removing a
debt adds'' rule --- the only place it is tested now, so it is worth checking. If a student writes $9-4$, $6-(-2)$ or $7-3$, that is the English
error, not an arithmetic one --- say so, and make them re-read the sentence.}

\qhead[cbred]{A4}{Difference}
\begin{enumerate}[label=\arabic*., itemsep=3pt]
  \item $8$
  \item $6$
  \item $4-(-13)=17$ degrees
  \item (a) $11$ \; (b) $-11$ \; (c) A difference is the gap between the numbers so it
        is never negative, but $-5-6$ is an instruction to move 6 further left.
  \item $18-(-24)=42$ degrees
\end{enumerate}

\qhead[cbred]{A5}{Word problems}
\begin{enumerate}[label=\arabic*., itemsep=3pt]
  \item $7-11+2=-2$. Floor $-2$ (B2).
  \item (a) $-3+11-14=-6\dC$. \; (b) Warmest $8\dC$ at noon, coldest $-6\dC$:
        $8-(-6)=14$ degrees.
  \item Still $-40$ dollars. Each week $+25-25=0$, so six weeks change nothing. This one
        rewards reading to the end before calculating.
\end{enumerate}

\qhead[cbred]{A6}{Now you write the question}
\noindent There is no single right answer --- check the four conditions in the green
box, in this order:
\begin{enumerate}[label=\arabic*., itemsep=3pt]
  \item \textbf{Does it start in the right place?} For (a) the story must \emph{begin}
        at $-5$, not arrive there. This is the same ``which number comes first''
        problem as A3, and it is the one most of them will get wrong.
  \item \textbf{Does the signal word match the sign?} (a) needs an \emph{up} word
        (rises, deposits, climbs) for the $+8$. (b) needs a \emph{removing a negative}
        --- a debt cancelled, a weight taken off --- not simply ``adds 4''.
  \item \textbf{Full sentences, ending in a question?}
  \item \textbf{Does their own answer come out as 3, and as 16?}
\end{enumerate}

\smallskip
\noindent{\small Two that work: (a) ``The temperature in Sa Pa was $-5\dC$ at
midnight. By noon it had risen by 8 degrees. What was the temperature at noon?''
(b) ``Mr Bowen had 12 dollars, and he also owed his brother 4 dollars. His brother told
him to forget the debt. How much is Mr Bowen worth now?'' \\
Part (b) is the harder one by a distance --- expect most of the class to write a plain
addition. That is worth a whole-class discussion rather than a red pen: ask them what
in real life makes you \emph{richer} when it is taken away.}

\qhead[cbred]{B1}{Match the word to its job}
\noindent 1--D \quad 2--F \quad 3--A \quad 4--C \quad 5--G \quad 6--H \quad 7--B \quad 8--E

\qhead[cbred]{B2}{Animal, plant, or both?}
\noindent Both: cell membrane, cytoplasm, nucleus, mitochondria. \\
Plant only: cell wall, chloroplast, sap vacuole. \\
{\small Nothing is animal-only. Expect some students to tick that column anyway ---
it is worth saying out loud that a plant cell has everything an animal cell has,
\emph{plus} extras.}

\qhead[cbred]{B3}{Answer in full sentences}
\begin{enumerate}[label=\arabic*., itemsep=3pt]
  \item A plant has no skeleton, so each cell needs a stiff wall around it to hold its
        shape and hold the plant up.
  \item Any two: the wall is thick and stiff while the membrane is thin and flexible;
        the wall is made of cellulose; the wall is only in plant cells while every cell
        has a membrane; the membrane controls what enters and leaves, the wall does not.
  \item They are chloroplasts, and the green substance inside them is chlorophyll.
  \item A microscope is a tool that uses curved pieces of glass called lenses to bend
        light and magnify an image, making it look much bigger than it really is.
\end{enumerate}

\qhead[cbred]{B4}{A model of a cell}
\noindent (a) bag --- cell membrane; jelly --- cytoplasm; plum --- nucleus;
beans --- mitochondria.

\smallskip
\noindent (b) Any one, with its organelle named: a shoebox or rigid container around the
bag --- cell wall; a water balloon in the middle --- sap vacuole; green sweets ---
chloroplasts.

\qhead[cbred]{B5}{Label the plant cell}
\noindent 1 --- cell wall \quad 2 --- cytoplasm \quad 3 --- cell membrane \quad
4 --- nucleus \quad 5 --- sap vacuole \quad 6 --- chloroplast \quad
7 --- mitochondrion

\fi

\end{document}
"

\documentclass[12pt,a4paper]{article}

% -- Answer key switch -------------------------------------------------------
\newif\ifkey
\ifdefined\TEACHER\keytrue\else
  \keyfalse        % \keyfalse = student copy   |   \keytrue = teacher copy
\fi

% -- House style -------------------------------------------------------------
% Shared by every packet. See homework/hw-style.tex.
% homework/hw-style.tex
% Shared house style for every Year 7 homework packet. \input this from the
% packet file, right after \documentclass and the answer-key switch.
%
% Do not fork this file per packet. If a packet needs a different look, the
% question is whether every packet should have it.
%
% The choices here are deliberate and were arrived at by printing and looking:
%   * No solid dark banners. Every header is a light tint with a coloured spine
%     on the left, so colour reads clearly without flooding a school printer.
%   * No marks and no mark totals. These are practice, not exam papers.
%   * Lato at 12pt with open leading, plus mathastext so inline maths is Lato
%     too. Without mathastext every "-6 + 4" is Computer Modern serif and the
%     sheet looks half-typeset.
%   * Writing lines are spaced for Year 7 handwriting, not adult margin notes.
%
% Provides: \wlines \blank \tickbox \sep \qmk-free marks (none) \hwsection
% \qhead \expr, the workspace box, and the remember / wordhelp / activity /
% yourturn coloured boxes. Palette matches the lesson decks exactly.

% -- Packages ----------------------------------------------------------------
\usepackage[T1]{fontenc}
\usepackage[utf8]{inputenc}
\usepackage[default]{lato}
\usepackage[a4paper,top=1.7cm,bottom=1.7cm,left=2.0cm,right=2.0cm,headheight=15pt]{geometry}
\usepackage{amsmath,amssymb}
\usepackage{xcolor}
\usepackage[most]{tcolorbox}
\usepackage{enumitem}
\usepackage{tikz}
\usetikzlibrary{arrows.meta}
\usepackage{array}
\usepackage{tabularx}
\usepackage{fancyhdr}
\usepackage{multicol}
\usepackage{needspace}
\usepackage{graphicx}

% Make maths use Lato too. Without this, every "-6 + 4" on the page is set in
% Computer Modern serif and the sheet looks half-typeset. Must come last.
\usepackage[italic]{mathastext}

\linespread{1.06}\selectfont

% -- The Learner's Book palette (same hex values as the lesson decks) ---------
\definecolor{cbteal}{HTML}{0087A8}
\definecolor{cbpurple}{HTML}{5C2483}
\definecolor{cborange}{HTML}{C25E12}
\definecolor{cbgreen}{HTML}{4A8B23}
\definecolor{cbred}{HTML}{C8102E}
\definecolor{cbblue}{HTML}{1A5FA8}
\definecolor{rulegray}{HTML}{9AA5AD}

% -- Answer lines and blanks -------------------------------------------------
% \wlines{n} draws n ruled writing lines, spaced wide enough for Year 7
% handwriting rather than for an adult with a fine pen.
\makeatletter
\newcommand{\wlines}[1]{%
  \par\vspace{0.75em}%
  \@tempcnta=0\relax
  \@whilenum\@tempcnta<#1\do{%
    \hbox to \linewidth{\color{rulegray}\leaders\hrule height 0.5pt\hfill}%
    \vspace{1.5em}%
    \advance\@tempcnta by 1\relax}%
  \par\vspace{-0.8em}%
}
\makeatother

% \blank[width] is an inline gap to write a short answer in.
\newcommand{\blank}[1][2.6cm]{\underline{\hspace{#1}}}

% A boxed space for showing working.
\newtcolorbox{workspace}[1][3.4cm]{%
  colback=white, colframe=rulegray, boxrule=0.5pt, arc=5pt,
  height=#1, left=6pt, right=6pt, top=4pt, bottom=4pt,
  before skip=9pt, after skip=11pt}

% An empty tick box.
\newcommand{\tickbox}{\fbox{\rule{0pt}{1.15em}\rule{1.15em}{0pt}}}

% A middot separator.
\newcommand{\sep}{\,\textperiodcentered\,}

% -- Section and question headings -------------------------------------------
% Light tint plus a fat coloured spine on the left. No solid dark banner: the
% colour reads clearly but barely costs the printer anything.
% needspace stops a heading stranding alone at the foot of a page. Kept modest:
% too large and it pushes headings over, leaving a worse hole than it prevents.
\newcommand{\hwsection}[2][cbteal]{%
  \needspace{8\baselineskip}%
  \par\vspace{13pt}%
  \noindent
  \begin{tcolorbox}[enhanced, nobeforeafter, width=\textwidth,
      colback=#1!7, colframe=#1!7, arc=5pt,
      borderline west={5pt}{0pt}{#1},
      left=14pt, right=10pt, top=7pt, bottom=7pt]
    {\large\bfseries\color{#1}#2}
  \end{tcolorbox}%
  \par\vspace{9pt}}

% \qhead[colour]{A2}{Moving along the line} -- a question heading.
\newcommand{\qhead}[3][cbteal]{%
  \needspace{6\baselineskip}%
  \par\vspace{11pt}%
  \noindent{\large\bfseries\color{#1}#2}\quad{\large\bfseries #3}\par\vspace{4pt}}

% -- Coloured boxes ----------------------------------------------------------
% Shared look: rounded, light fill, coloured rule, coloured title text on a
% pale title strip rather than a dark one.
\tcbset{hwbox/.style={
  enhanced, breakable, arc=6pt, boxrule=1pt,
  left=11pt, right=11pt, top=8pt, bottom=8pt,
  fonttitle=\bfseries, attach boxed title to top left={xshift=12pt, yshift=-8pt},
  boxed title style={arc=4pt, boxrule=1pt, left=7pt, right=7pt, top=2pt, bottom=2pt},
  top=13pt}}

% Orange: things to remember. Matches the deck's "Write This Down".
\newtcolorbox{remember}[1][Remember from the lesson]{hwbox,
  colback=cborange!6, colframe=cborange, coltitle=cborange,
  colbacktitle=white, title={#1}}

% Blue: English help. The barrier in this class is the wording, not the maths.
\newtcolorbox{wordhelp}[1][Word help]{hwbox,
  colback=cbblue!6, colframe=cbblue, coltitle=cbblue,
  colbacktitle=white, title={#1}}

% Purple: an activity or instruction.
\newtcolorbox{activity}[1][How to do this homework]{hwbox,
  colback=cbpurple!6, colframe=cbpurple, coltitle=cbpurple,
  colbacktitle=white, title={#1}}

% Green: your turn / a friendly nudge.
\newtcolorbox{yourturn}[1][Your turn]{hwbox,
  colback=cbgreen!7, colframe=cbgreen, coltitle=cbgreen!75!black,
  colbacktitle=white, title={#1}}

% -- Shorthands --------------------------------------------------------------
\newcommand{\dC}{\,^{\circ}\mathrm{C}}          % use inside maths: $-8\dC$
\newcommand{\key}[1]{{\color{cborange}\textbf{#1}}}
\newcommand{\eng}[1]{{\color{cbblue}\textbf{#1}}}

\setlist[enumerate]{itemsep=8pt, topsep=5pt, leftmargin=*}
\setlist[itemize]{itemsep=4pt, topsep=4pt, leftmargin=*}
\renewcommand{\arraystretch}{1.45}

% -- An expression the student is asked to work from -------------------------
\newcommand{\expr}[1]{%
  \tcbox[on line, colback=cbgreen!10, colframe=cbgreen!55, boxrule=1pt, arc=4pt,
         left=9pt, right=9pt, top=4pt, bottom=4pt]{\large\bfseries $#1$}}

% -- Running head and foot ---------------------------------------------------
% The packet overrides \fancyhead[L] with its own title.
\pagestyle{fancy}
\fancyhf{}
\fancyhead[L]{\small\color{black!55}Year 7\sep Homework}
\fancyhead[R]{\small\color{black!55}Name: \blank[3.4cm]}
\fancyfoot[C]{\small\color{black!55}Page \thepage}
\renewcommand{\headrulewidth}{0pt}
\renewcommand{\headrule}{\vspace{2pt}\hbox to\headwidth{\color{cbteal!45}\leaders\hrule height 1.2pt\hfill}}



% -- Per-packet running head -------------------------------------------------
\fancyhead[L]{\small\color{black!55}Year 7\sep Homework 1}

\begin{document}

% ===========================================================================
% COVER BLOCK
% ===========================================================================
\thispagestyle{fancy}

\begin{tcolorbox}[enhanced, colback=cbpurple!7, colframe=cbpurple!7, arc=7pt,
                  borderline west={6pt}{0pt}{cbpurple},
                  left=18pt, right=14pt, top=14pt, bottom=14pt]
  {\color{cbpurple}\bfseries YEAR 7\sep UNIT 1\sep HOMEWORK 1}\\[6pt]
  {\color{cbpurple}\Huge\bfseries Integers and Cells}\\[10pt]
  {\color{black!70}Mathematics 1.1 --- Adding \& Subtracting Integers}\\
  {\color{black!70}Science 1.1 --- Cells and Microscopes}
\end{tcolorbox}

\vspace{10pt}
\noindent
\begin{tabularx}{\textwidth}{@{}X X@{}}
  \textbf{Name:} \blank[5.6cm] & \textbf{Class:} \blank[5.6cm] \\
  \textbf{Date given:} \blank[4.6cm] & \textbf{Date due:} \blank[5.0cm] \\
\end{tabularx}

\vspace{8pt}
\begin{activity}
  \begin{itemize}
    \item Write your answers \textbf{on this sheet}, in the spaces given. If you need
          more room, use the back of the page.
    \item In maths, \key{write the calculation before the answer}.
    \item In science, answer in \key{full English sentences}. ``Chloroplast'' is not a
          sentence; ``The green circles are chloroplasts'' is.
    \item If a word is new to you, look for the blue \eng{Word help} box on that page
          before you give up.
  \end{itemize}
\end{activity}

% ===========================================================================
% SECTION A -- MATHEMATICS
% ===========================================================================
\hwsection{Section A \textperiodcentered\ Mathematics 1.1 --- Adding and Subtracting Integers}

\begin{remember}
  An \key{integer} is a whole number that is positive, negative or zero --- never a
  fraction. On the number line, \key{negative} integers sit to the \key{left} of zero
  and \key{positive} integers to the \key{right}.

  \medskip
  \begin{tabularx}{\linewidth}{@{}lX@{}}
    Add a \key{positive}      & move \key{right} \\
    Add a \key{negative}      & move \key{left} \\
    Subtract a \key{positive} & move \key{left} \\
    Subtract a \key{negative} & move \key{right} --- because $a-(-b)=a+b$ \\
  \end{tabularx}

  \medskip
  Careful: only $-(-)$ becomes $+$. A single $+(-)$ still sends you left, so
  $5+(-3)=2$, \textbf{not} $8$.
\end{remember}

% -- A1 ---------------------------------------------------------------------
\qhead{A1}{Show it on the number line}

\noindent\textbf{(a)} Show $4 + (-3)$ on the number line below. Put a cross on the
number you \textbf{start} at, then draw an arrow for the move.

\vspace{12pt}
\noindent
\resizebox{0.92\textwidth}{!}{%
\begin{tikzpicture}[x=1cm,y=1cm]
  \draw[line width=1.1pt,cbteal,{Latex[length=3mm]}-{Latex[length=3mm]}] (-5.7,0) -- (8.7,0);
  \foreach \n in {-5,...,8}{
    \draw[line width=0.9pt,cbteal] (\n,-0.18) -- (\n,0.18);
    \node[below=4pt,font=\small] at (\n,-0.18) {$\n$};
  }
\end{tikzpicture}}

\vspace{10pt}
\noindent\hspace*{0.5em}$4 + (-3) = \blank[2.2cm]$

\vspace{12pt}
\noindent\textbf{(b)} Complete each sentence with \eng{higher than}, \eng{lower than},
\eng{warmer than} or \eng{colder than}.

\begin{enumerate}[label=(\roman*)]
  \item $-8$ is \blank[3.4cm] $-4$.
  \item A temperature of $3\dC$ is \blank[3.4cm] a temperature of $-9\dC$.
  \item Floor $-2$ of the car park is \blank[3.4cm] floor $1$.
\end{enumerate}

% -- A2 ---------------------------------------------------------------------
\qhead{A2}{Moving along the line}

\noindent Work these out. You may draw a small number line beside any of them.

\vspace{10pt}
\begin{multicols}{3}
\begin{enumerate}[label=(\alph*), itemsep=13pt]
  \item $-6+4 = \blank[1.6cm]$
  \item $5+(-9) = \blank[1.6cm]$
  \item $-3-5 = \blank[1.6cm]$
  \item $-2-(-7) = \blank[1.6cm]$
  \item $8+(-8) = \blank[1.6cm]$
  \item $-10-(-4) = \blank[1.6cm]$
  \item $0-(-6) = \blank[1.6cm]$
  \item $-7+15 = \blank[1.6cm]$
  \item $4-11 = \blank[1.6cm]$
  \item $-5-(-5) = \blank[1.6cm]$
  \item $6-8-(-5) = \blank[1.6cm]$
  \item $-9-6+(-1) = \blank[1.6cm]$
\end{enumerate}
\end{multicols}

% -- A3 ---------------------------------------------------------------------
\qhead{A3}{Say it, then write it}

\begin{wordhelp}[Word help --- these words decide the order of your calculation]
  \begin{tabularx}{\linewidth}{@{}lX@{}}
    \eng{subtract 5 from 8}   & you \textbf{start} at 8: $8-5$. The English gives the numbers in the \textbf{opposite order} to the calculation. \\
    \eng{take 6 away from 2}  & the same trap: $2-6$. \\
    \eng{7 less than 3}       & again the same trap: $3-7$. \\
    \eng{rises by / falls by} & the temperature goes up / the temperature goes down. \\
    \eng{deposit}             & to put money \textbf{into} an account $(+)$. \\
    \eng{withdraw}            & to take money \textbf{out of} an account $(-)$. \\
    \eng{owe} / \eng{a debt}  & money you must pay back. It makes your total \textbf{negative}. \\
    \eng{descend}             & to go \textbf{down}. To go up is to \eng{ascend}. \\
  \end{tabularx}
\end{wordhelp}

\noindent Turn each English sentence into a calculation, then work out the answer.
\key{Write the calculation first.}

\vspace{6pt}
\begin{enumerate}[label=\arabic*., itemsep=11pt]
  \item Subtract 9 from 4. \\[4pt]
        Calculation: \blank[6.2cm] \quad Answer: \blank[2.4cm]
  \item Take 6 away from $-2$. \\[4pt]
        Calculation: \blank[6.2cm] \quad Answer: \blank[2.4cm]
  \item What is 7 less than 3? \\[4pt]
        Calculation: \blank[6.2cm] \quad Answer: \blank[2.4cm]
  \item The temperature is $-8\dC$ and it rises by 5 degrees. \\[4pt]
        Calculation: \blank[6.2cm] \quad Answer: \blank[2.4cm]
  \item A submarine is at $-30$ m. It descends a further 12 m. \\[4pt]
        Calculation: \blank[6.2cm] \quad Answer: \blank[2.4cm]
  \item An account holds $-35$ dollars. The bank removes a debt of 20 dollars from it.
        \\[4pt]
        Calculation: \blank[6.2cm] \quad Answer: \blank[2.4cm]
\end{enumerate}

% -- A4 ---------------------------------------------------------------------
\qhead{A4}{Difference}

\begin{remember}[Remember: a difference is a gap]
  The \key{difference} between two numbers is how far apart they are on the line.
  Work it out with \key{bigger $-$ smaller}. A difference is \key{never negative}.
  \eng{How much warmer}, \eng{how much colder}, \eng{how much higher} and
  \eng{how many more} all ask for the same thing.
\end{remember}

\begin{enumerate}[label=\arabic*.]
  \item Find the difference between $-6$ and $2$. \hfill Answer: \blank[2.8cm]
  \item Find the difference between $-9$ and $-3$. \hfill Answer: \blank[2.8cm]
  \item How much warmer is $4\dC$ than $-13\dC$? Show your calculation.
        \wlines{1}
  \item These two questions look almost the same. They are not.
        \begin{enumerate}[label=(\alph*), itemsep=9pt]
          \item Find the difference between $-5$ and $6$. \hfill Answer: \blank[2.8cm]
          \item What is $-5$ minus $6$? \hfill Answer: \blank[2.8cm]
          \item In one full sentence, explain why the two answers are different.
                \wlines{2}
        \end{enumerate}
  \item One winter morning it is $18\dC$ in Ho Chi Minh City and $-24\dC$ in
        Ulaanbaatar. How much colder is Ulaanbaatar? Show your calculation.
        \wlines{1}
\end{enumerate}

% -- A5 ---------------------------------------------------------------------
\qhead{A5}{Word problems}

\noindent Read each one \key{all the way to the end} before you start writing.
Show the calculation, then the answer.

\begin{enumerate}[label=\textbf{\arabic*.}, itemsep=16pt]

  \item \textbf{The office lift.} Mr Bowen gets into the lift on floor 7. He goes
        \textbf{down} 11 floors, and then \textbf{up} 2 floors. Which floor is he on
        now?
        \begin{workspace}[2.8cm]\end{workspace}

  \item \textbf{A day in Sa Pa.} At 6 a.m. the temperature is $-3\dC$. By noon it has
        risen 11 degrees. By midnight it has fallen 14 degrees.
        \begin{enumerate}[label=(\alph*)]
          \item What is the temperature at midnight?
          \item What is the \textbf{difference} between the warmest and the coldest
                temperature of the day?
        \end{enumerate}
        \begin{workspace}[3.4cm]\end{workspace}

  \item \textbf{Six weeks of saving.} Mr Bowen's account is at $-40$ dollars. Every
        Monday he deposits 25 dollars. Every Friday he withdraws 25 dollars. How much is
        in the account after 6 weeks? Explain your answer in one full sentence.
        \begin{workspace}[3.1cm]\end{workspace}

\end{enumerate}

% -- A6 ---------------------------------------------------------------------
\qhead[cbgreen]{A6}{Now you write the question}

\begin{yourturn}[Your turn to be the teacher]
  So far the questions have given you the English, and you have written the maths.
  Now do it \textbf{backwards}. You are given the maths --- write the English.

  \medskip
  Your word problem must:
  \begin{itemize}
    \item be about something \textbf{real} --- money, temperature, floors in a
          building, or depth under the sea;
    \item use at least one word from the \eng{Word help} box on the earlier page,
          such as \eng{rises by}, \eng{falls by}, \eng{deposits}, \eng{withdraws},
          \eng{owes}, \eng{below} or \eng{descends};
    \item be written in \textbf{full English sentences}, and end with a question;
    \item have the answer that is already given to you.
  \end{itemize}
\end{yourturn}

\noindent\textbf{(a)} Write a word problem for \quad \expr{-5 + 8 = 3}
\wlines{4}

\vspace{14pt}
\noindent\textbf{(b)} Write a word problem for \quad \expr{12 - (-4) = 16}
\wlines{4}

% ===========================================================================
% SECTION B -- SCIENCE
% ===========================================================================
\hwsection{Section B \textperiodcentered\ Science 1.1 --- Cells and Microscopes}

\begin{remember}
  A \key{cell} is the smallest basic unit of all living organisms. An \key{organelle}
  is a tiny structure inside a cell that does one specific, important job to keep the
  cell alive. Every cell has a \key{cell membrane}, \key{cytoplasm}, a \key{nucleus}
  and \key{mitochondria}. A plant cell has all of those \key{plus} a \key{cell wall}
  made of \key{cellulose}, \key{chloroplasts} containing \key{chlorophyll}, and a
  \key{sap vacuole}.
\end{remember}

% -- B1 ---------------------------------------------------------------------
\qhead{B1}{Match the word to its job}

\noindent Write the correct letter in the box next to each word.

\vspace{8pt}
\noindent
\begin{tabularx}{\textwidth}{@{}p{4.2cm} @{\hspace{6pt}} c @{\hspace{12pt}} X@{}}
  \textbf{Word} & \textbf{Letter} & \textbf{Job} \\[2pt]
  1.\ Cell            & \tickbox & \textbf{A} \; The control centre of the cell --- the ``boss''. \\
  2.\ Organelle       & \tickbox & \textbf{B} \; A large space of sugars and water that keeps a plant cell firm. \\
  3.\ Nucleus         & \tickbox & \textbf{C} \; Where energy is released from food. \\
  4.\ Mitochondria    & \tickbox & \textbf{D} \; The smallest basic unit of all living organisms. \\
  5.\ Cell wall       & \tickbox & \textbf{E} \; A tool that uses lenses to bend light and magnify an image. \\
  6.\ Chloroplast     & \tickbox & \textbf{F} \; A tiny structure inside a cell that does one specific job. \\
  7.\ Sap vacuole     & \tickbox & \textbf{G} \; A strong, stiff outer layer that holds a plant cell in shape. \\
  8.\ Microscope      & \tickbox & \textbf{H} \; Where a plant makes its food, using sunlight. \\
\end{tabularx}

% -- B2 ---------------------------------------------------------------------
\qhead{B2}{Animal, plant, or both?}

\noindent Put a tick ($\checkmark$) in \textbf{one} column for each organelle.

\vspace{8pt}
\noindent
\begin{tabularx}{\textwidth}{|X|>{\centering\arraybackslash}p{3.1cm}|>{\centering\arraybackslash}p{2.9cm}|>{\centering\arraybackslash}p{1.9cm}|}
  \hline
  \textbf{Organelle} & \textbf{Animal cells only} & \textbf{Plant cells only} & \textbf{Both} \\
  \hline
  Cell membrane & & & \\ \hline
  Cytoplasm     & & & \\ \hline
  Nucleus       & & & \\ \hline
  Mitochondria  & & & \\ \hline
  Cell wall     & & & \\ \hline
  Chloroplast   & & & \\ \hline
  Sap vacuole   & & & \\ \hline
\end{tabularx}

% -- B4 ---------------------------------------------------------------------
\qhead{B3}{Answer in full sentences}

\begin{wordhelp}[Sentence starters --- you may copy these]
  ``A plant cell needs a cell wall because\ldots''\\
  ``One difference is that\ldots\ Another difference is that\ldots''\\
  ``A microscope is a tool that\ldots''
\end{wordhelp}

\begin{enumerate}[label=\arabic*., itemsep=14pt]

  \item An animal has a skeleton to hold it up. A plant does not. Explain why every
        plant cell needs a cell wall.
        \wlines{2}

  \item Give \textbf{two} ways a cell wall is different from a cell membrane.
        Write each one as a full sentence.
        \wlines{4}

  \item Under the microscope, a leaf cell is full of small green circles. Name them,
        and name the green substance inside them.
        \wlines{2}

  \item What does a microscope do? Use the words \key{lens} and \key{magnify} in your
        answer.
        \wlines{3}

\end{enumerate}

% -- B5 ---------------------------------------------------------------------
\qhead{B4}{A model of a cell}

\noindent Mr Bowen builds a model of an animal cell in his kitchen. He uses a clear
plastic bag, filled with jelly, with a plum and four red beans pushed into the middle
of it.

\vspace{8pt}
\noindent\textbf{(a)} Which organelle does each item represent?

\vspace{6pt}
\noindent
\begin{tabularx}{\textwidth}{|p{5cm}|X|}
  \hline
  \textbf{Item in the model} & \textbf{Organelle it represents} \\
  \hline
  The plastic bag \rule{0pt}{1.9em} & \\ \hline
  The jelly       \rule{0pt}{1.9em} & \\ \hline
  The plum        \rule{0pt}{1.9em} & \\ \hline
  The four beans  \rule{0pt}{1.9em} & \\ \hline
\end{tabularx}

\vspace{14pt}
\noindent\textbf{(b)} Name \textbf{one} thing he could add to turn it into a model of a
\textbf{plant} cell, and say which organelle it would represent.
\wlines{2}

% -- B3 ---------------------------------------------------------------------
\qhead{B5}{Label the plant cell}

\noindent Write the name of each numbered part on the lines below the diagram.

\vspace{12pt}
\noindent
\resizebox{\textwidth}{!}{%
\begin{tikzpicture}[x=1cm,y=1cm,
    lbl/.style={circle, draw=cbteal, fill=white, line width=1.1pt,
                inner sep=1.8pt, minimum size=8mm, font=\bfseries}]

  % -- the cell itself. The wall band is deliberately 0.7 wide so that the
  %    wall (1) and the membrane (3) are visibly two different things.
  \fill[cbgreen!35, rounded corners=12pt] (0,0) rectangle (11,6.4);        % cell wall
  \fill[cbteal!10, rounded corners=8pt]  (0.7,0.7) rectangle (10.3,5.7);   % cytoplasm
  \draw[line width=1.2pt, black!60, rounded corners=8pt] (0.7,0.7) rectangle (10.3,5.7);
  \draw[line width=1.6pt, cbgreen!70!black, rounded corners=12pt] (0,0) rectangle (11,6.4);

  % -- sap vacuole --
  \fill[cbblue!18] (5.5,2.95) ellipse (3.2 and 1.65);
  \draw[line width=0.9pt, cbblue!60] (5.5,2.95) ellipse (3.2 and 1.65);

  % -- nucleus --
  \fill[cbpurple!30] (2.0,4.35) circle (0.85);
  \draw[line width=0.9pt, cbpurple!80] (2.0,4.35) circle (0.85);
  \fill[cbpurple!60] (2.0,4.35) circle (0.28);

  % -- chloroplasts --
  \foreach \x/\y/\r in {1.7/1.5/20, 3.2/1.15/-15, 5.0/1.05/10,
                        6.8/1.15/-20, 8.5/1.5/15, 9.6/2.6/-10, 9.6/4.2/25}{
    \begin{scope}[shift={(\x,\y)}, rotate=\r]
      \fill[cbgreen!75] (0,0) ellipse (0.46 and 0.24);
      \draw[line width=0.7pt, cbgreen!45!black] (0,0) ellipse (0.46 and 0.24);
    \end{scope}}

  % -- mitochondria. Kept away from x = 5.5 so the cytoplasm leader (2) has a
  %    clear, empty patch of cytoplasm to land on.
  \foreach \x/\y/\r in {3.3/5.1/-12, 8.4/5.1/16, 1.35/2.7/70}{
    \begin{scope}[shift={(\x,\y)}, rotate=\r]
      \fill[cborange!65] (0,0) ellipse (0.42 and 0.2);
      \draw[line width=0.7pt, cborange!80!black] (0,0) ellipse (0.42 and 0.2);
      \draw[line width=0.5pt, cborange!80!black] (-0.25,0.05) -- (-0.05,-0.06) -- (0.15,0.06);
    \end{scope}}

  % -- leader lines and numbers. Each one ends on a dot so there is no doubt
  %    about which structure it is pointing at.
  \foreach \sx/\sy/\ex/\ey/\lx/\ly/\n in {%
      2.2/6.4/2.2/7.55/2.2/7.95/1,          % cell wall  (the green band)
      5.5/5.35/7.9/7.55/8.0/7.95/2,         % cytoplasm  (empty patch)
      0.7/4.9/-1.5/6.25/-1.9/6.4/3,         % cell membrane (the black line)
      2.0/4.35/-1.5/4.05/-1.9/4.0/4,        % nucleus
      5.5/2.95/12.6/3.3/13.0/3.35/5,        % sap vacuole
      1.7/1.5/-1.5/1.1/-1.9/1.05/6,         % chloroplast
      8.4/5.1/12.6/5.5/13.0/5.55/7}{        % mitochondrion
    \draw[line width=0.7pt, black!65] (\sx,\sy) -- (\ex,\ey);
    \fill[black!65] (\sx,\sy) circle (0.075);
    \node[lbl] at (\lx,\ly) {\n};}

  % keep the bounding box sensible
  \path (-2.6,-0.5) rectangle (13.9,8.6);
\end{tikzpicture}}

\vspace{16pt}

% A tabular rather than multicols, so the seven answer lines can never be split
% across a page break away from their diagram.
\noindent
\begin{tabular}{@{}p{0.48\textwidth} p{0.48\textwidth}@{}}
  1.\ \blank[5.4cm] & 5.\ \blank[5.4cm] \\[30pt]
  2.\ \blank[5.4cm] & 6.\ \blank[5.4cm] \\[30pt]
  3.\ \blank[5.4cm] & 7.\ \blank[5.4cm] \\[30pt]
  4.\ \blank[5.4cm] &                   \\
\end{tabular}

\vspace{10pt}
\noindent{\small\color{black!60}Hint: number 3 is the very thin layer just inside the
outer green edge. Numbers 1 and 3 are easy to confuse --- one is thick and stiff, one is
thin and flexible.}

% ===========================================================================
% ANSWER KEY (teacher copy only)
% ===========================================================================
\ifkey
\newpage
\hwsection[cbred]{Answer Key --- Teacher Copy}

% The key is for an adult reading it at a desk, not a 12-year-old writing on it,
% so tighten everything back up.
\linespread{1.0}\selectfont
\setlist[enumerate]{itemsep=2pt, topsep=3pt, leftmargin=*}
\setlist[itemize]{itemsep=2pt, topsep=3pt, leftmargin=*}

\qhead[cbred]{A1}{Show it on the number line}
\noindent (a) Cross on $4$, arrow three steps \emph{left}, landing on $1$. So
$4+(-3)=1$. The cross must be on 4, not on 3 --- adding a negative means starting at
the first number and moving left.\\
(b) (i) lower than \quad (ii) warmer than \quad (iii) lower than.
{\small\ (``higher than'' is also acceptable in (iii) if the sentence is reversed ---
mark the reasoning, not the word.)}

\qhead[cbred]{A2}{Moving along the line}
\noindent
(a) $-2$ \quad (b) $-4$ \quad (c) $-8$ \quad (d) $5$ \quad (e) $0$ \quad (f) $-6$ \\
(g) $6$ \quad (h) $8$ \quad (i) $-7$ \quad (j) $0$ \quad (k) $3$ \quad (l) $-16$

\smallskip
\noindent{\small Common errors: (d) and (f) --- students flip \emph{any} pair of signs
to $+$. Only $-(-)$ becomes $+$. (b) $5+(-9)$ is a single $+(-)$, so it still goes
left.}

\qhead[cbred]{A3}{Say it, then write it}
\begin{enumerate}[label=\arabic*., itemsep=3pt]
  \item $4-9=-5$
  \item $-2-6=-8$
  \item $3-7=-4$
  \item $-8+5=-3$
  \item $-30-12=-42$
  \item $-35-(-20)=-15$
\end{enumerate}
\noindent{\small Questions 1--3 are the ``reversed order'' trap; 6 is the ``removing a
debt adds'' rule --- the only place it is tested now, so it is worth checking. If a student writes $9-4$, $6-(-2)$ or $7-3$, that is the English
error, not an arithmetic one --- say so, and make them re-read the sentence.}

\qhead[cbred]{A4}{Difference}
\begin{enumerate}[label=\arabic*., itemsep=3pt]
  \item $8$
  \item $6$
  \item $4-(-13)=17$ degrees
  \item (a) $11$ \; (b) $-11$ \; (c) A difference is the gap between the numbers so it
        is never negative, but $-5-6$ is an instruction to move 6 further left.
  \item $18-(-24)=42$ degrees
\end{enumerate}

\qhead[cbred]{A5}{Word problems}
\begin{enumerate}[label=\arabic*., itemsep=3pt]
  \item $7-11+2=-2$. Floor $-2$ (B2).
  \item (a) $-3+11-14=-6\dC$. \; (b) Warmest $8\dC$ at noon, coldest $-6\dC$:
        $8-(-6)=14$ degrees.
  \item Still $-40$ dollars. Each week $+25-25=0$, so six weeks change nothing. This one
        rewards reading to the end before calculating.
\end{enumerate}

\qhead[cbred]{A6}{Now you write the question}
\noindent There is no single right answer --- check the four conditions in the green
box, in this order:
\begin{enumerate}[label=\arabic*., itemsep=3pt]
  \item \textbf{Does it start in the right place?} For (a) the story must \emph{begin}
        at $-5$, not arrive there. This is the same ``which number comes first''
        problem as A3, and it is the one most of them will get wrong.
  \item \textbf{Does the signal word match the sign?} (a) needs an \emph{up} word
        (rises, deposits, climbs) for the $+8$. (b) needs a \emph{removing a negative}
        --- a debt cancelled, a weight taken off --- not simply ``adds 4''.
  \item \textbf{Full sentences, ending in a question?}
  \item \textbf{Does their own answer come out as 3, and as 16?}
\end{enumerate}

\smallskip
\noindent{\small Two that work: (a) ``The temperature in Sa Pa was $-5\dC$ at
midnight. By noon it had risen by 8 degrees. What was the temperature at noon?''
(b) ``Mr Bowen had 12 dollars, and he also owed his brother 4 dollars. His brother told
him to forget the debt. How much is Mr Bowen worth now?'' \\
Part (b) is the harder one by a distance --- expect most of the class to write a plain
addition. That is worth a whole-class discussion rather than a red pen: ask them what
in real life makes you \emph{richer} when it is taken away.}

\qhead[cbred]{B1}{Match the word to its job}
\noindent 1--D \quad 2--F \quad 3--A \quad 4--C \quad 5--G \quad 6--H \quad 7--B \quad 8--E

\qhead[cbred]{B2}{Animal, plant, or both?}
\noindent Both: cell membrane, cytoplasm, nucleus, mitochondria. \\
Plant only: cell wall, chloroplast, sap vacuole. \\
{\small Nothing is animal-only. Expect some students to tick that column anyway ---
it is worth saying out loud that a plant cell has everything an animal cell has,
\emph{plus} extras.}

\qhead[cbred]{B3}{Answer in full sentences}
\begin{enumerate}[label=\arabic*., itemsep=3pt]
  \item A plant has no skeleton, so each cell needs a stiff wall around it to hold its
        shape and hold the plant up.
  \item Any two: the wall is thick and stiff while the membrane is thin and flexible;
        the wall is made of cellulose; the wall is only in plant cells while every cell
        has a membrane; the membrane controls what enters and leaves, the wall does not.
  \item They are chloroplasts, and the green substance inside them is chlorophyll.
  \item A microscope is a tool that uses curved pieces of glass called lenses to bend
        light and magnify an image, making it look much bigger than it really is.
\end{enumerate}

\qhead[cbred]{B4}{A model of a cell}
\noindent (a) bag --- cell membrane; jelly --- cytoplasm; plum --- nucleus;
beans --- mitochondria.

\smallskip
\noindent (b) Any one, with its organelle named: a shoebox or rigid container around the
bag --- cell wall; a water balloon in the middle --- sap vacuole; green sweets ---
chloroplasts.

\qhead[cbred]{B5}{Label the plant cell}
\noindent 1 --- cell wall \quad 2 --- cytoplasm \quad 3 --- cell membrane \quad
4 --- nucleus \quad 5 --- sap vacuole \quad 6 --- chloroplast \quad
7 --- mitochondrion

\fi

\end{document}
